
\vspace{2em}

次日，一纸堂堂正正、却字字如刀的檄文，猝然划破了内域维持许久的冷战僵局。

檄文由叶牧谦一人所撰，以云隐别院之名发出；以刀削斧凿的文字，以其理念为核心，向整个内域发出雷霆般的号召：所有认同“问道途”之念、苦于守旧联盟压迫与盘剥已久的修士与凡民，当共举义旗，涤荡寰宇，争一个公理，求一条活路！

檄文在内域各处被飞速地传抄，其文辞激越昂扬，理直气壮，更因蕴含着一股受尽屈辱后迸发出的决绝意志，而具有了灼人的力量。

其中两句，尤为耀眼夺目，如同烧红的烙铁，狠狠地烫在了每一个读到它的人心头上：

“我未失礼于彼，彼乃无礼横行，专恃兵坚器利，自取决裂如此！”

“与其苟且图存，贻羞万古，孰若大张挞伐，一决雌雄？”

檄文精准地道出了所有被压迫、被孤立、被践踏者心中积郁已久、几乎要喷薄而出的愤懑与那份置之死地而后生的决绝。

\vspace{2em}

在那座象征着内域传统权力巅峰的紫阳宗凌霄殿内，副宗主、守旧联盟此战的实际统帅陈炎，同样手持一份抄本。

他端坐于主位之上，指尖漫不经心地轻轻敲打着玉简，脸上竟也浮现出一丝复杂的、近乎玩味的笑意，低声自语道：“好骂。”

但这声赞叹里，浸透了居高临下的冰冷轻蔑与毫不掩饰的嘲讽。

大张挞伐？一决雌雄？

陈炎几乎要嗤笑出声。

在他看来，这种宣战甚至完全是有利于他们一方的：以前或许还有所顾忌、不敢轻易越雷池一步；但战书一下，双方自然可以堂堂正正、短兵相接了。

檄文越激愤，在陈炎的眼中，便越显得可笑而悲凉，如同困兽最后的哀鸣。

他的自信源于冷酷到极致的现实计算：

紫阳宗有着久经操练、传承有序的精锐弟子，他们演练纯熟的战阵攻防一体；丹鼎阁有着庞大的资源储备仓库，其门下擅长辅助与诡异毒术的药修、毒修被大量编入作战序列；而已经被刘长靖彻底掌控的无极剑派及其众多附属势力，则能够提供最为锋利、擅长攻坚与突袭的剑锋。

而一个崛起不过数载、根基浅薄、全靠一群被排挤的散修和少数几个叛逆宗门弟子支撑起来的新生势力，其底蕴拿什么来抗衡紫阳宗、丹鼎阁、无极剑派这三大巨头数百上千年积累的庞大资源、深厚功法和无数强横修士？

就算己方不主动进攻，仅凭旷日持久的资源消耗和硬实力的绝对碾压，也足以将云霞山上那点可怜的反抗力量慢慢磨死、困死。

\vspace{2em}

云霞山上。

虽然说叶牧谦和吴东浩形成了共识，不能再继续绥靖以紫阳宗为首的老大势力的联合绞杀，必须要重拳回击，他们在诸弟子前往青鸾秘境的过程中，在外界也确实拉起了一支由散修构成的义军；但在如何回击上，双方确实产生了分歧。

此刻，云隐别院之内，关于具体战术的实施，正在进行一场争论。

与其说是争论，不如说是一场理念与情绪的直接对撞。

青石板周围聚集了三四十人，除了叶牧谦、白言冰、陈千羽、沈瑶，以及很早就跟随叶牧谦随修的陈默等人，其余皆是吴东浩凭借自身威望与吴刚的情报网网罗来或主动汇聚来的、自愿组成义军的散修头领或代表人物。

洞府内的空气因灵灯与众人身上散发的驳杂气息，显得有些浑浊。

叶牧谦本人虽然没有完整地读过《论持久战》，但是其精髓倒是早早地刻进了他的心中。

而他的观点自然也是非常朴实无华的：“‘问道途’有民心汇聚之功，此战长远来看，我方必胜；但按如今实力来看，敌强我弱，我义军不能正面撄其锋芒，故不能速胜。叶某认为，短期来看，义军主力必须固守少数据点；其余人等化整为零，潜入广大乡村地区，对敌军游而击之，并广泛传播我问道途理念，团结更多群众于我旗帜之下，陷敌人于汪洋大海，然后胜之。”

吴东浩还没发话，草图另一侧，一位一身劲装、身形挺拔如枪的另一位散修头目立即插言。

头目名为尹壮图，修为正在玄脉境通络期。

他声音洪亮，斩钉截铁：“叶先生的方略，深远周全，尹某佩服。可我等举义，聚拢人心，靠的不是画地为牢、被动挨打！这么多年了，兄弟们从四面八方赶来，心里头都憋着一股火呢！”

他猛地一挥手，目光灼灼地扫过在场许多面露赞同、甚至隐含焦躁的散修头领：“我们必须要一场胜仗——一场堂堂正正、能让所有人都看见的胜仗！一场能把陈炎那老匹夫的嚣张气焰打下去的胜仗！提气，壮胆，更要告诉全内域的兄弟姐妹，我们不仅敢立旗，更敢拔剑，能砍下那群老爷们的脑袋！”

他的话语如同火星溅入干草堆，瞬间点燃了许多人的情绪。

这些散修大多是在旧有秩序下备受排挤、挣扎求存的个体，能修炼至今，靠的便是一股子狠劲和不择手段的拼搏；话语更是直接、血性，充满了草莽豪气，远比叶牧谦那些迂回而忍耐的策略更合胃口。

“尹兄说得在理！”

“对！咱爷们儿聚在这儿，不是来当缩头乌龟的！”

“总得干他娘的一票大的，让紫阳宗知道疼！”

“是啊，叶先生，咱们的实力未必就不如他们一股偏师，找个机会，狠狠咬下一块肉来！”

附和声此起彼伏，不少人的眼睛亮了起来，仿佛已经看到了伏击成功、凯旋而归的场景。

那种渴望用胜利证明自己、用敌人鲜血洗刷憋屈的心情，几乎要凝成实质。

叶牧谦静静地听着，脸上依旧平静，只是眉宇间掠过一丝几不可察的凝重。

待声浪稍平，他才缓缓开口，声音一如既往的平稳，却带着一种穿透喧嚣的清晰：“尹兄所言，是为提振士气，凝聚人心，此心可鉴。然，战争非意气之争，更非江湖械斗之放大。”

他向前半步，手指再次落回青石板，沿着代表敌军推进方向的几道箭头移动。“敌军大致分三路而来，看似有隙可乘。陈炎用兵，看似笨重迟缓，实则步步为营，稳扎稳打，其战阵互为犄角，灵力探测法阵与巡逻警戒网交错密布，几乎没有留给大规模伏击部队隐蔽接近、从容设伏的空隙。”

他顿了顿，看向诸多散修，“即便我们能集结一支足以撼动其一路偏师的精锐力量，成功突袭，最快能多久结束战斗？半个时辰？一个时辰？”

尹壮图立即被噎住了。

吴东浩则思考了一下：虽然说支持义军的声势确实浩大，但是这么长时间以来，真正上山、愿意拼命作战的不过两万左右，而修为稍高、到达玄脉境界者区区数千。

但数千玄脉境，再怎么说也已经不是一个小数目了。

吴东浩一估计，得出了结论：“若谋划得当，突袭迅猛，速战速决，一个时辰内当可击溃其一部前锋！”

“一个时辰。”叶牧谦重复道，指尖在敌阵后方一点，“足够另外两路，甚至中军主力做出反应。届时，我方伏击部队大概率陷入夹击或苦战；即便伏击取胜，缴获些许物资，杀伤数百敌玄脉境界及以上的修士，于敌军整体而言不过九牛一毛。然我军可能付出的代价，将是数百甚至上千最精锐、最敢战之士的折损。”

他抬起头，目光扫过那些情绪激昂的头领：“以我有限之血勇精锐，赌敌军无穷之资源消耗；况且，败则伤筋动骨，胜亦无济于事，此实是下下之选！”

另一位散修头目，叫曹润生的，脸色有些涨红，争辩道：“叶先生！打仗哪有不冒险的？若是瞻前顾后，何时能打开局面？我们一味退守，只会让他们越发骄狂！兄弟们心里这口气不顺，修为都难有寸进！我并非要浪战，而是寻机歼敌一部，哪怕两败俱伤，也要打出我‘问道途’的血性！让天下人看看，我们不是任人拿捏的软柿子！”

叶牧谦自己也有些忍不住了：“你们这是胡闹，用好不容易聚集起来的主力去胡闹！”

他再次指向草图上那片代表敌人腹地、统治薄弱的地域，“决胜之力，在彼不在我。过早暴露我们有限的主力精锐，将其消耗在正面碰撞中，乃是舍本逐末。”

不同的主张清晰地摆在了所有人面前：诸多散修要的，是一场立竿见影、提振人心的正面战斗，哪怕代价惨重；叶牧谦坚持的，则是隐忍固守，积蓄力量，将胜负手放在更深远、也更需要时间的敌后发动与群众争取上。

会场的氛围变得微妙而凝重。散修们大多面露不甘，他们觉得叶牧谦太过书生气，太过保守，甚至有些怯战。

在他们混迹底层的经验里，“百无一用是书生”，那种书生气是要不得的，必须靠着一股不要命的狠劲杀出条血路才行！

退让和隐忍往往意味着失去机会，随即人为刀俎、我为鱼肉！

叶牧谦描绘的“星火燎原”固然美好，但太过遥远：远水难解近渴。

眼下，他们需要的是胜利，是证明，是宣泄！

而叶牧谦这边，除了白、陈、沈三位无条件相信他的判断之外，吴刚也拧着眉头，看看正气势汹汹的尹壮图、曹润生等，又看看叶牧谦，似乎内心也有些挣扎。

他理解那些散修的躁动，因为那也曾是他的心态——他变得长袖善舞，那已经是后话了。

眼见得会场内部隐隐有了分裂的迹象，吴东浩深吸一口气，知道言语上的辩论已难见效。这时候，便是他发挥能力的时候了。

和稀泥！

他转向众人，抱拳道：“诸位兄弟，是非曲直，难以论说。我等既举义旗，便是军中同袍。军中最重号令统一，而号令出自众议。我提议，就此战术方略进行表决。在场诸位，皆有表决之权。无论结果如何，吴某必当遵从！”

此言一出，附和声再次响起，且更加响亮。对于这些散修而言，这种“公投”式的方式，远比听从某一人——尤其是叶牧谦这样“不通实务”的学术领袖——的决断，更让他们感到被尊重。

叶牧谦眼帘微垂，心中无声地叹息。

这就是人心！

在巨大而切近的压力下，在渴望即时反馈的情绪驱使下，更具煽动性、更符合直接经验的主张，往往能压倒更理性、更深远但见效慢的计划。

他也没有反对表决，因为强行压制只会离心离德。

“行。便依吴兄所言。”

表决的过程很快，几乎毫无悬念。

白言冰、陈千羽、沈瑶，自然站在叶牧谦一边；

吴刚、吴东浩弃权。

而其余三十余名散修头领或代表，超过九成，齐刷刷地举手，或者出声，明确表示支持“主动寻机出击”的方略。他们的眼神热烈，带着对胜利的渴望。

相比之下，叶牧谦那套守正待机、“决胜于千里之外”的理论，在他们听来有些缥缈，远不如“带兄弟们砍翻他们”来得实在。更何况，叶牧谦温文尔雅，更像一位令人尊敬的师长、学者；而尹壮图、曹润生则浑身散发着刀头舔血、敢打敢拼的“自己人”气息。

结果显而易见。

尹壮图再次抱拳，声音铿锵：“承蒙诸位兄弟信任！尹某在此立誓，必当竭尽所能，寻隙而击，为咱们‘问道途’义军，打出第一场漂亮的胜仗！”

欢呼声在洞府内响起，充满了对即将到来战斗的期待。

叶牧谦安静地站在原地，青衫素雅，与周围热血沸腾的氛围显得有些疏离。白言冰担忧地看向师尊，陈千羽抿了抿唇，沈瑶则目光微冷地扫过那些欢呼的散修。

待声浪稍歇，叶牧谦才缓缓开口，声音依旧平静无波：“既然众议已决，那……吴兄便为义军总指挥，节制诸人；尹壮图、曹润生则为副手，总揽寻机出战事宜。”他看向吴东浩，又扫了一眼尹、曹二人，目光深邃，“望你们审时度势，慎选战机，保全我义军元气为要。”

尹、曹二人郑重抱拳：“叶先生放心，我等省得！”

叶牧谦点了点头，不再多言。

打一打，也好。

实践是检验真理的唯一标准，即使实践存在代价。而搞一场会战，也确实能真正摸一摸敌人的底细。

当然，在叶牧谦本人看来，这场会战必然失败。他相信吴东浩的镇场能力，只希望尹、曹二人能别把义军主力糟蹋得元气大伤了。

\vspace{2em}

数日之后，义军开拔，一万余人浩浩荡荡地下了山，本来因为“山林啸聚”一番而鼎沸的云霞山顿时冷清了不少，只剩下少数后勤部。

叶牧谦、白言冰、陈千羽、沈瑶均未随军。

望着那七拼八凑但气冲斗牛的、渐行渐远的义军部队，叶牧谦深深地叹了一口气。

“瑶儿，走吧，再去加固一下九宫迷踪阵。言冰，你去协助那些后勤部队，在敌人可能攻上来的地方埋设陷阱、建设工事。千羽，你去看看药田，保证能够按时收获。”

四人各个心乱如麻。

“未来大概率会有一场守山的苦战了。”叶牧谦叹气道。

\vspace{2em}

这是一片位于云霞山以北约三百里、两片低矮丘陵夹峙下的宽阔浅谷。

按照尹壮图与几位心腹部下反复侦察推演得出的结论，这里是守旧联盟西路一支偏师——由紫阳宗某位外门长老孙翊率领、约八千人的战营——前往预定集结地的必经之路，且因其地形相对闭塞（虽然也没有那么闭塞），大军难以迅速展开，是打一场迅猛伏击的理想场所。

吴东浩本人倒也谨慎，一万多人的义军，留下了两千士兵在更远的地方试图接应、作为预备队。

前来埋伏的，大约万人。

尹壮图站在沟壑一侧的隐蔽高点上，劲装被山风吹得紧贴身躯，猎猎作响。

他目光如鹰隼，扫视着下方逐渐进入伏击圈的敌军队伍。那支队伍军容严整，士卒身着赤红的制式灵甲，行动间隐隐有灵力共鸣的微光流转，步伐虽然因山路而稍显缓滞，但前后呼应严密，斥候游骑交替前出，显示出良好的训练素养。

他身边，曹润生半蹲着，手中一枚符箓微微闪光，低声道：“尹兄，看旗号，是紫阳宗的一部，主将是那个叫孙翊的长老，听说修为在玄脉后期。他们队列中央那些覆盖着毡布的大车，灵力波动不弱，怕是运送的补给或某种法器。”

尹壮图舔了舔有些干裂的嘴唇，眼中非但没有惧色，反而燃起更炽烈的战意。

“紫阳宗的赤皮狗，装备倒是光鲜。正好，宰了这条大鱼，扒下他们的皮，给兄弟们换换装！告诉各队，按原计划，等他们前军过了一半，中军那些大车进入沟底最窄处，听我号令，一起动手！先集中火力轰击车队，打乱他们阵脚，分割前后，然后近身绞杀！”

在他们的规划下，中军将是主要的突破口，在这里安排了六千人之多。而在突袭发起之后，袭扰前军和后军的，分别不过千人。

命令被悄无声息地传递下去。埋伏在沟壑两侧乱石、灌木后的上万义军修士，此刻屏息凝神。他们大多眼神狂热，紧握着自己五花八门的兵器——有的只是凡铁打造、附加了基础锋锐符文的刀剑，有的是品相不一的低阶飞剑，更有拿着奇门钩锁、淬毒短弩的。

虽然装备驳杂，但那股子被压抑了许久的、渴望用敌人鲜血证明自己的杀意，却几乎凝成实质，让空气都变得粘稠而躁动。

吴东浩隐在另一侧的山崖后，眉头紧锁。

他同意了这次出击，但内心深处那丝不安始终挥之不去。

他看得比尹壮图更清楚：敌人数量虽略比我方少，但其阵型严谨，反应速度绝不会慢。

义军人数虽多，可一旦第一波突袭不能达到预期效果，陷入缠斗，后果不堪设想。

“盯紧侧翼和后方，一旦情况有变，准备接应尹壮图他们撤出。”

此刻，那支队伍，如同一条赤红色的巨蟒，缓缓游入黑云沟的腹腔。阳光被两侧山崖遮挡，沟内光线昏暗，只有甲胄和兵刃偶尔反射出冰冷的幽光。队伍中隐约传来军官的呵斥和车轮碾过碎石的声响。

前军过了谷口，辎重正好进入沟底最深处。

时机到了！

尹壮图猛地站起，灵力毫无保留地爆发开来。他手中那柄门板似的厚重巨剑高高举起，剑身嗡鸣，亮起土黄色的沉重光芒。

“兄弟们！随我杀——！”

“杀啊——！！！”

“为了问道途！为了活路！”

积蓄已久的怒吼声，如同山洪暴发，从沟壑两侧轰然炸响！刹那间，成千上万道灵力光芒亮起，虽然驳杂不纯，却汇聚成一片令人目眩的死亡潮汐，朝着沟底的敌军倾泻而下！

最先落下的是各种远程攻击。火球、冰锥、风刃、雷符、毒镖、飞针……密密麻麻，如同夏日狂暴的冰雹，劈头盖脸砸向尚未完全反应过来的紫阳宗队伍。

尤其是针对那十几辆覆盖毡布的大车，更是受到了重点关照，超过百道攻击汇聚成数股洪流，狠狠撞了上去！

轰！轰轰轰——！

剧烈的爆炸声接连响起，火光冲天，浓烟翻滚。几辆大车瞬间被炸得粉碎，里面的灵石、矿石、甚至尚未组装完成的弩炮零件四处飞溅。

拉车的驯化妖兽发出凄厉的哀嚎，受惊之下横冲直撞，进一步搅乱了敌军队列。

“敌袭！御！”

紫阳宗队伍中响起了尖锐的警哨和军官声嘶力竭的吼叫。

训练有素的精锐在这一刻展现了价值，尽管遭遇突袭，出现了一些伤亡和混乱，但士卒仍然在极短时间内做出了反应。赤红色的灵光从他们甲胄上升腾而起，迅速连成一片，形成一面巨大的、半弧形的火焰护盾，将后续的大部分攻击抵挡在外。

盾牌铿锵组合，长矛如林竖起，瞬间构筑起一道坚固的防线。

但义军的突袭势头太猛，第一波打击也确实造成了可观杀伤。近百名紫阳宗修士在最初的混乱中非死即伤，队伍被拦腰截断，前后失去了联系。

“就是现在！冲下去！近战！杀光他们！”尹壮图身先士卒，如同一头发狂的猛虎，从数丈高的崖壁上一跃而下，巨剑带着开山裂石之势，狠狠斩向一名刚刚组织起小队、试图反击的紫阳宗什长。

那什长举盾格挡，但尹壮图含怒一击何等狂暴，只听“咔嚓”一声脆响，精铁盾牌连同其后的一条手臂应声而断，鲜血喷溅！

尹壮图去势不减，巨剑横扫，又将旁边两名惊愕的紫阳宗修士拦腰斩飞！

“尹首领威武！”

“杀啊！”

主将的悍勇极大鼓舞了士气，无数义军修士红着眼睛，从埋伏点冲下，如同决堤的洪水，涌向被分割开的敌军段落；灵枢境的找灵枢境的大头兵，玄脉境的找玄脉境的长官，捉对厮杀。

短兵相接，血肉横飞的惨烈战斗，瞬间在黑云沟狭窄的底部上演。

个人武艺与悍不畏死的精神，在这一刻似乎取得了优势。

许多散修出身、在底层摸爬滚打练就了一身厮杀本领的义军，三五成群，凭借对地形的熟悉和一股不要命的狠劲，围着那些试图结阵自保的小股紫阳宗修士猛攻。

刀光剑影，符篆乱飞，惨叫与怒吼混杂在一起，浓烈的血腥味迅速弥漫开来。

曹润生也杀入了战团，他使得是一对分水峨眉刺，身形滑溜如鱼，专挑敌人防护薄弱处下手，出手狠辣刁钻，几个照面便放倒了两名紫阳宗士卒。他看到战局似乎正向有利于己方的方向发展，脸上忍不住露出一丝兴奋的潮红，高喊道：“兄弟们加把劲！吃掉他们！”

分配给敌人前军和后军的小股部队也开始袭扰起来，迟滞着敌人的回援，短期内成果斐然。

就连远处观战的吴东浩，看到初期战果，紧绷的心弦也略微松了一分。

难道……叶先生的判断还真是过于保守了？

我们真能啃下这块硬骨头？

然而，他们高兴得太早了。

孙翊此刻已从最初阵型被搅乱的惊恐中镇定下来，冷哼一声，命令通过灵力传遍己方尚能联系的部队。

“传令！中军抛弃辎重，自行结阵固守！前后两部，起焚城大阵，左右卫护，缓步压上，与中军靠拢！”

想了想，又对近卫补充道：

“发求援焰火，告知西路周圣子：我部遇伏，敌军近万，请求合围！”

命令迅速被执行。

只见那些尚未被义军彻底缠住的中军紫阳宗修士，迅速放弃与冲下来的散修纠缠，开始向后收缩。他们彼此靠拢，大约每百名修士便迅速聚集成一个紧密的圆阵，外围盾牌层层叠叠，长矛从缝隙中刺出，内里的修士则开始整齐划一地掐动法诀，赤红色的火行灵力在他们之间疯狂汇聚、共鸣。

“不好！他们要结战阵！”一名有见识的义军小头目骇然变色，大声示警。

但已经晚了。

中军的一个个防御圆阵已经结成，变成铁桶，打不动了！

更要命的是，前军和后军已经结成了更具有攻击性的战阵！

“烈焰焚城，焚尽八荒！疾！”

孙翊暴喝一声，手中令旗猛地挥下！

嗡——！

众多紫阳宗修士组成的战阵中心，空气剧烈扭曲，一团团直径超过十丈、炽白到让人无法直视的巨大火球骤然成型，散发出恐怖的高温，连周围的岩石都开始融化、流淌！

下一刻，这团毁灭性的火球轰然炸开，化作无数道手臂粗细、灵动如蛇的炽白火线，以战阵为中心，呈放射状向着四面八方疯狂蔓延、抽打、绞杀！

这正是紫阳宗核心战阵之一的威力！

嗤嗤嗤——！

火线所过之处，无论是义军修士仓促撑起的护体灵光，还是他们手中的低劣法器，甚至是坚韧的肉身，都如同热刀切猪油般被轻易撕裂、洞穿、点燃！

“啊——！我的腿！”

“火！扑不灭！”

“退！快退开！”

惨叫声此起彼伏，原本气势如虹、扑到近前的义军，顿时遭受重创。

数十名冲在最前面的悍勇之士，瞬间被火线穿透，变成一个个燃烧的火人，凄厉地翻滚着，很快化为一堆焦炭。更多的人被火线擦中，护体灵光剧烈闪烁后破碎，身上燃起难以扑灭的灵火，惨叫着向后逃窜。

原本还算有序的进攻浪潮，被这突如其来的、覆盖性的恐怖打击彻底打懵、打散！

而这，仅仅是个开始。

孙翊显然意图将这股胆大包天的伏兵全部留下。

战阵在第一波齐射后，开始如同一个个燃烧的刺猬，缓缓向前、向沟内压迫移动。火线不断从战阵中衍生、抽击，清扫着前方一切障碍。

同时，战阵内部的修士轮换施法，一道道范围稍小但同样致命的大片法术被抛射出来，落在义军人群相对密集的区域。

黑云沟内，顿时化为了火焰地狱。原本昏暗的沟底被映照得一片通红，炽热的气浪扭曲升腾，将空气都烧得滚烫。焦臭味、皮肉烧灼的可怕气味混合着血腥，形成了令人作呕的死亡气息。

义军最大的劣势——缺乏统一指挥和应对集团战阵的经验——在此刻暴露无遗。

他们习惯了小规模的混战和游击，也对此颇有自得；何曾见过如此恐怖、如此有组织性的灵力覆盖打击？

慌乱之下，众多小队队长的命令也混乱起来：

“顶上去！”

“跑啊，等死呢！”

“我听哪个！”

命令互相冲突，导致局部更加混乱。

有小队想从侧翼绕过战阵，但那战阵移动间毫无破绽，很快全队死的死、逃的逃。

有人想远程攻击打断施法，可零星的法术打在厚重的联合灵力护盾上，连涟漪都难以激起。

尹壮图目眦欲裂！

他亲眼看着自己麾下几名得力干将，在试图组织一波反冲锋时，被交叉扫过的七八道炽白火线直接气化，连渣都没剩下。

他狂吼一声，将巨剑舞得密不透风，磕飞了几道袭向自己的火线。虽然尹壮图本人修为较高、扛下了众多低阶修士组成的超大型战阵的攻击；但那火线蕴含的灼热灵力和冲击力，震得他气血翻腾，手臂发麻。

“不能退！退了就全完了！跟我上，杀了孙翊那家伙！！”

尹壮图知道，不击破中军、不停止孙翊的动作，今日所有人都要葬身火海。

尹壮图与曹润生集结了身边还能指挥的约两百名最精锐、最不怕死的修士，鼓起残存的灵力，化作一股决死的尖刀，试图强行凿穿战阵外围的防御，强行换死孙翊，停止他对毁灭性超大战阵的指挥。

这是一场鸡蛋碰石头的自杀式冲锋。

“螳臂当车！”孙翊在阵中看得分明，嘴角泛起残忍的冷笑。

“剑修！给我撕了他们！”

早已在战阵侧后方蓄势待发的五百名无极剑派外剑脉剑修，如同一群闻到血腥味的鲨鱼骤然杀出！他们的任务，就是应对这种敌方精锐的突围或反击。

这些剑修身法极快，行动间带着森然的统一性。他们没有结什么复杂战阵，只是简单分为数支小队，如同几柄淬毒的匕首，倏忽之间便切入尹壮图组织的冲锋队列侧翼和薄弱衔接处。

剑光！

冰冷、迅疾、精准、高效的剑光！

他们的剑法远远不如已被清理的心剑脉那般灵动缥缈、直指本心，但在杀伐效率上，却达到了某种极致。

每一剑都直奔要害，配合默契，往往两三人一组，瞬间便能绞杀一名冲在前面的义军好手。

“呃啊！”曹润生一个不慎，左腿被一道刁钻的剑气划过，深可见骨，鲜血狂喷。

他闷哼一声，险些栽倒，被旁边一名修士拼死拖住。

尹壮图的冲锋势头被硬生生遏止，陷入了与这些精锐剑修的死斗之中。他巨剑势大力沉，每一击都足以开碑裂石，但那些外剑修士滑溜异常，根本不与他硬拼，只是游斗、袭扰、分割，将他身边稍不留神的战友们一个个快速点杀。

他空有一身勇力，却仿佛陷入了一张无形而坚韧的剑网，憋屈得几欲吐血。

就在尹壮图部陷入苦战、整个义军战线因为核心战阵的碾压和外剑修士的切割而开始动摇、溃散的危急关头——

呜——！

低沉雄浑的号角声，从黑云沟的两端入口处，同时响起！声音苍凉悠远，却带着令人心悸的肃杀之意。

紧接着，地平线上，烟尘大作！旌旗招展，灵力冲霄！

西路讨伐军的真正主力，到了！

而且不止一路！从沟壑的东、西两个出口方向，同时出现了滚滚而来的钢铁洪流。东面是更多的紫阳宗战阵，赤旗如血；西面则是丹鼎阁与部分附属宗门的联军，旗帜花色繁杂，但阵型同样严整。他们显然接到了孙翊的求援，并做出了最迅速、也最致命的反应——反包围。

周文彬本人没出现，似乎是算准了用不着他出手。

原本计划打一场速战速决伏击战的义军，此刻反而成了瓮中之鳖。

“完了……”吴东浩远远看到这一幕，一颗心沉到了谷底。

最坏的情况发生了。

敌人的反应速度和兵力调动效率，远超他们的预估。陈炎“稳扎稳打，互为犄角”的策略确实是将强大的实力化为了难以撼动的整体优势。

他们这支孤军深入的义军主力，已经成了送上门的肥肉！

“吴首领！东、西两面都有大队敌军出现！我们被包围了！”浑身浴血的斥候连滚爬爬地冲过来报信，脸上满是绝望。

吴东浩手青筋暴起，他猛地看向沟底仍在拼死搏杀的尹壮图部，又看看周围已经开始出现恐慌、骚动的义军部队。他知道，必须立刻做出决断，能撤出去多少是多少！

“传令！各部不要恋战！以小队为单位，自行突围！能走多少是多少！尹首领那边……我去接应！”吴东浩咬牙。

他知道凭自己真符境摹形期的修为，在许多个成熟战阵中救出尹壮图希望渺茫；但作为义军总指挥，他不能眼睁睁看着主将战死而不救，那会让士气彻底崩溃。

然而，战场局势的崩坏速度，比吴东浩预想的还要快。

当两路生力军加入战场，并且迅速在外围构建起更严密的封锁线，开始用远程法器和战阵技能向沟内进行覆盖式轰击时，义军的崩溃便不可避免地发生了。

丹鼎阁的修士隐匿在侧翼，将大片大片的甜腻毒瘴借助风势吹入沟内，许多修为较低的义军吸入后立刻手脚酸软，灵力运转停滞，成了待宰羔羊；地面在土行法术的干扰下变得泥泞不堪，甚至突然塌陷，严重阻碍了突围的脚步。

绝望像瘟疫一样蔓延。原本凭借血勇支撑的斗志，在绝对的实力差距和绝境面前，迅速消融。

有人扔下武器，跪地投降，但换来的往往是毫不留情的一剑。

更多人像没头苍蝇一样乱撞，然后被各色法术光芒吞噬。

尹壮图身边的战友越来越少，最后只剩寥寥十余人，被超过百名的紫阳宗精锐和外剑修士团团围住在一个小山坡上。

他浑身是伤，巨剑拄地，大口喘着粗气，眼神却依旧桀骜凶狠，死死盯着不远处被重重保护的周焱。

“贼子！可敢与尹某单挑！”他嘶声吼道，声音沙哑。

孙翊嗤笑一声，甚至懒得回答，只是摆了摆手。

顿时，更加密集的法术和剑气向那小山坡倾泻而去。

“保护首领！”最后几名亲卫怒吼着扑上，用身体挡住了大部分攻击，瞬间被淹没在光芒中。

尹壮图狂吼，压榨出最后一丝灵力，巨剑爆发出最后的黄光，一式力劈华山，斩碎了袭向自己的几道剑气和火球。但他已是强弩之末，护体灵光黯淡如风中残烛。

噗嗤！噗嗤！

两柄飞剑，一左一右，如同毒蛇吐信，趁机突破了他巨剑的防御范围，狠狠贯入了他的胸膛和腹部！

尹壮图高大的身躯猛地一震，他低头看了看透体而出的剑尖，脸上闪过一丝茫然，随即化为无尽的愤怒与不甘。

“悔……不该……”他喃喃吐出几个字，眼神的光芒迅速黯淡下去，最终轰然倒地，激起一片尘土。那双曾经燃烧着炽热战意的眼睛，至死未曾闭合，直直望着云霞山的方向。

几乎在尹壮图倒下的同时，吴东浩才杀入重围。

见到尹壮图已经陨落，也只能长啸一声，扛起尸体再次强行杀出。

真符境毕竟不是法相境，虽然说这么多战阵留不下他，但要说硬碰硬，那还是远远不够的。

另一边，曹润生早在腿受重创时就被亲信架着后撤，但突围路上同样遭遇了数次截杀，身上又添了几处新伤，最重的一处在后背，是被一道散射的剑气所伤，差点伤及灵枢，人也因失血过多而昏迷过去。

主将战死，副将重伤。吴东浩此刻留下的后手——在远处接应的数千修士——确实起到了自己的战术作用，成功发动了突袭、搅乱了些战场局势，但已经不能改变溃败的结果了。

而原来埋伏敌人的、剩下的修士彻底失去了组织，只能凭借本能亡命奔逃。

那已经不能称之为撤退，而是一场彻头彻尾的溃败。

守旧联盟的军队则如同经验丰富的猎手，开始有条不紊地追杀、分割、歼灭溃逃的残敌。

一开始参与埋伏的近万人，折了五千有余，剩余的拼死拼活才逃了出来；预备队伤亡倒是不大，但也心有余悸。

不少人本因为自己被编入预备队而感到不忿，但此刻只剩下更大的庆幸。

夕阳如血，缓缓沉入西边的山峦，将这片修罗场映照得一片凄艳。

喊杀声渐渐平息，取而代之的是伤者压抑的哀嚎、胜利者打扫战场的喧嚣，以及食腐妖兽被血腥味吸引而来的低沉嘶吼。

曾经气势汹汹而来的一万两千余义军修士，最终能够狼狈逃出的，不足八千。而且绝大多数带伤，建制彻底打散，士气低落到了极点。

更惨重的损失是高层：义军骁将尹壮图确认战死；曹润生重伤昏迷，脊背的伤势可能留下永久性的隐患，即便醒来，战力也大不如前。吴东浩本人修为高绝，在背负尹壮图尸体回来的途中受了些轻伤，所幸无大碍。

其余在战斗中表现出色、有一定威望的中下层头目，伤亡更是超过三分之一。

经此一役，云霞山义军可谓元气大伤，不仅损失了大量最敢战、最有经验的骨干，更重要的是，那口下山时“欲与天公试比高”的心气，被彻底打没了。

恐惧、沮丧、质疑、乃至对叶牧谦当初劝阻的后悔情绪，开始在残存者中悄然蔓延。

而对守旧联盟而言，此战役则是一场预料之中的胜利。虽然己方阵亡修士也大约二千有余，但比例上看并不算多；更重要的是，它以极小的代价重创了对手的有生力量，还验证了己方战略战术的正确性，极大地鼓舞了士气。

陈炎在接到战报后，只是淡淡地评价了四个字：“乌合之众。”

通往云霞山的最后屏障，正在这场血腥的大败后，缓缓洞开。真正的铁壁合围与攻坚，即将开始。而云霞山上，叶牧谦站在反复加固过的阵眼枢纽处，接到染血的败报时，脸上并无多少意外之色，只有一片深沉的凝重与决绝。

实践确实检验了真理。只是这代价，未免太过沉重。

随后，在战役失败的总结与反思会议上，叶牧谦自然而然地接过了义军的总指挥权。

叶牧谦不禁喟叹：每次他接手的总是一片烂摊子！

但斯大林格勒从来不相信眼泪。

\vspace{2em}

大规模会战宣告失败后的一年内时间，便在无数场规模不等、却同样浴血的小规模接战，在一次次带着不甘与悲愤的撤退，在一声声无奈放弃经营许久据点的叹息中，迅速流逝。

吃下一地，稳固一地，再吃下下一地！

一年，原问道途所有地下情报与走私网络，近乎全毁！

没毁的也不敢再出来活动了！

守旧联盟的大军，如同一个巨大而沉重的石磨，缓慢、笨重，却带着无可抗拒的力量，一寸一寸地碾过了问道途势力范围的外围城镇、资源点和战略据点。

最终，这碾压一切的兵锋，携带着胜利者的傲慢与肃杀，直抵那最后的、也是最高的屏障——云霞山。

这一日，陈炎亲率中军主力，旌旗蔽空，终于兵临云霞山下。

站在山巅俯瞰，景象足以让最坚韧的战士也感到一阵窒息般的绝望。

山下，守旧联盟的营寨依着地势，连绵铺展开去，一眼望不到边际，仿佛给大地铺上了一层铁灰色的毡毯。

各色代表不同宗门、不同战营的旌旗如林般矗立，在风中猎猎作响，几乎遮天蔽日。

从那些规划严整的营寨中，升腾起一道道、一片片色泽各异的灵力光晕。

赤红的火行灵力灼热逼人，青木的生机中暗藏毒瘴，白金剑气锋锐刺目，土黄厚重如山，水蓝绵密森寒……

这些光晕交织、融合在半空，形成了一片瑰丽奇诡却充满纯粹杀伐之气的灵力光云，冲霄而上，将云霞山半边天空都渲染得变了颜色，连日光都显得黯淡昏沉。

出入山林的每一条大小要道，都被鹿砦、壕沟以及时刻运转的警戒与探测法阵几乎彻底锁死，巡逻修士的队伍甲胄鲜明，气息精悍，交错往复，目光如鹰隼般扫视着山上的动静。

山下，是阵列森严的钢铁森林，是汹涌澎湃、近乎无限的灵力海洋，是志在必得、以泰山压顶之势而来的碾压洪流。山上，是历经一年苦战、人人面带疲惫却眼神依旧不屈的战士，是日益见底的丹药瓶与灵石箱，是寄托了内域无数受压迫者最后希望与信仰的险要壁垒。

一副瓮中捉鳖的恶劣态势，已然在战场层面彻底形成。

战争的初级阶段，以守旧联盟凭借绝对硬实力的无情碾压，和问道途一方被迫陷入艰苦卓绝的战略防御、乃至被重重围困的绝境而告终。

那赤裸裸的、令人绝望的实力差距，在这场从一开始就极不对等的生死较量中，展现得淋漓尽致，冰冷而残酷。

然而，正是在这山下连营的肃杀之气几乎要凝结成实体、如同铅云般沉沉压上云霞山巅、让人喘不过气来的时刻，坐镇于云隐别院深处、那掌控整座山防御体系核心阵眼处的叶牧谦，其神色却异乎寻常地平静，甚至可以说得上是淡漠。

他独自立于阵眼枢纽的灵光之中，青衫素雅，身形挺拔如孤峰。

那双仿佛能洞悉世间一切虚妄与表象的眼睛，平静地望向山下那密密麻麻的营寨，望向那冲霄而起、炫耀着武力的斑斓灵力光云。

眸底深处，并无半分寻常人应有的慌乱、恐惧或是愤怒，只有一片如同古井深潭般的、深邃的了然。

过去两年多看似平静实则暗流汹涌的冷战岁月里，过去一年正面战场被不断推进、绞杀的日子里，除了指挥义军撤退，他倾注最多心血、智慧与珍贵资源的，便是此云霞山本身了。

他有信心，九宫迷踪能够将脚下这座云霞山，从一座普通的险峻山峰，经营成一座纵使面对千军万马、拥有雄厚资源的敌人，也无法被轻易碾碎、啃下的战争堡垒。

而作为整个堡垒核心的，自然就是他的得意之作：九宫迷踪阵！

“传令各部，”叶牧谦的声音通过嵌入山体灵脉的特定传讯法符，清晰、稳定、不带丝毫波澜地回荡在山上各处要害据点守备负责人的耳中。

“放弃外围所有隘口、前出哨所，全部撤回云霞山主峰预设防御圈内。”

他略一停顿，那平静的声音却下达了一个将决定接下来无数人生死的命令：

“即刻起，开启九宫迷踪阵第一重变。”

命令如山，早已在无数次演练中刻入骨髓的守军，此刻如同上好了发条的精密器械，开始沉默而高效地运转。

散布在山脚下各处险要地点，原本准备誓死阻击、拖延时间的修士小队，毫不犹豫地放弃了经营许久的阵地。

埋设预设的阻挠陷阱后，小队们沿着早已勘测熟稔的隐秘小径，迅速而有序地向主峰核心区域收缩。

在敌我实力差距悬殊到令人绝望之时，主动将伸出的五指攥紧成拳，收回胸膛，看似退缩，实则是为了将力量凝聚于一点。

将那看似柔软的棉絮层层包裹内部最坚硬的铁芯，准备在敌人最意想不到、也最能发挥己方优势的位置，给予其最凶狠、最致命的反击。

这，是叶牧谦前世阅读那些波澜壮阔的历史时，从一次次以弱胜强、于绝境中开辟新生的经典战例中，领悟到的精髓。

叶牧谦自认无法像前世的伟人一样有着极高的军事天赋与才能，而现在的条件也不允许：正面兵力不足以打出如雷霆之势摧枯拉朽的三大战役，他自己的指挥能力也不足以打出四渡赤水这样精彩且精确的战役。

但如果仅仅是打一个反围剿，那似乎容易不少，至少他有很多现成的答案可以抄。

\vspace{2em}

云霞山，本就以山势险峻、路径错综复杂著称。

强行仰攻本就是兵家大忌，而加上九宫迷踪阵，对攻山的守旧联盟而言，情况只会更差。

就在守旧联盟的先锋部队，以严整战阵为依托，沿着几条主要山道气势汹汹、志在必得地向上推进之时，原本清晰可辨的山路、熟悉的溪涧与巨石，仿佛被一只无形而庞大的手轻轻搅动、抹去，然后重新排列。

浓密得化不开、伸手不见五指的白雾，不知从岩缝中、树根下还是地脉里汹涌而出，顷刻间如同涨潮的乳白色海水，弥漫、吞没了整个山林外围区域。雾气中蕴含着被阵法强行扭转的紊乱灵气，如同无数细密的针，疯狂刺扰着闯入者的神识探查：

斥候刚刚回头，信誓旦旦地报告前方路径正常，转瞬间，身后的主力部队就骇然发现，方才走过的坚实山路已然消失，脚下竟是云雾缭绕的陌生断崖；

明明按照记忆和地图该向左急转的路口，踏出去的队伍却直通一个布满苔藓湿滑、且犬牙交错的天然石刺绝地。整片外围山区，在短短半盏茶的时间内，便化作了一座庞大、诡异、时刻变化的活迷宫。

这，还仅仅是令人头痛的序曲。

九宫迷踪阵的扭曲空间效应，与预先巧妙埋设在关键节点的无数陷阱相结合，产生了叠加的恐怖效果：

当敌军队列在浓雾中如盲人般艰难摸索时，脚下看似坚实的落叶覆盖的地面可能突然无声塌陷，露出下方深达数丈、底部插满淬毒精钢刺的陷坑，瞬间吞噬数名躲闪不及的修士；

头顶上方那仿佛亘古存在的岩壁会毫无征兆地整体松动，数十根被符文加持过、裹挟着万钧之力的合抱滚木与棱角狰狞的礌石轰然砸落，破风之声如同死神的咆哮；

更隐蔽且防不胜防的，是那些触发式的小型连环灵力爆符，或是被特殊法术催生的、见血便疯狂滋长的铁线荆棘藤蔓种子。

这些陷阱虽难以对结成战阵、有层层护身灵光笼罩的高阶修士造成瞬间致命伤，却足以在关键时刻打乱他们严谨的推进节奏，迫使整个战阵停滞或变形，更在不断消耗他们宝贵的灵力以维持防御，并在每一名敌军心中持续施加无形的压力，制造挥之不去的混乱与随时丧命的恐慌。

守旧联盟想象中摧枯拉朽的推进速度，仅仅在山脚之下，就被这层黏稠、危险、无所不在的迷雾沼泽极大地迟滞、拖住了。

\vspace{2em}

天师府。

作为内域名义上地位超然的中立势力，其总坛不断向冲突双方发出正式文函，呼吁“秉持天道好生之德，即刻停战，以和谈化解干戈”。

紫阳宗为首的守旧联盟自然对此嗤之以鼻，势在必得；云霞山一方更无妥协余地。

但天师府的举动，远不止于这些流于表面的公文往来。

总坛深处，幽静的观星阁内，府主沈天穹一袭朴素道袍，正凭栏远眺，目光似乎穿过了千山万水，落在了烽火连天的云霞山方向。他的指尖无意识地在冰冷的玉石栏杆上轻轻敲击，那节奏缓慢而深沉，如同他此刻反复权衡的内心。

对叶牧谦此人，沈天穹的情感颇为复杂。自女儿沈瑶拜师，以及后来那次开诚布公的长谈后，他便一直在关注这个看似年轻之人的动向。

叶牧谦所倡的“问道途”，那种打破陈规、务实求真、愿为更多修士乃至凡人开辟前路的理念，在沈天穹看来，固然激进，甚至有些理想主义的天真；但其内核的精神，却与他年轻时所怀的某种抱负隐隐共鸣。

那是一种不愿被旧有秩序完全束缚、渴望看见新变化的微妙欣赏。

然而，身处他这个位置，个人的好恶必须让位于全局的权衡。

沈天穹收回远眺的目光，转身望向阁内悬挂的内域堪舆巨图：图上写写画画，已经很旧了。

云霞山的位置被一枚红色的符文标记着，而代表守旧联盟兵锋的黑色箭头，正从四面八方向其合围。

他看得十分透彻：陈炎所率联军，实力占据绝对上风，若任由其速战速决，以雷霆万钧之势碾碎云霞山，那么此后内域的格局将彻底失衡。

一个迅速膨胀、吞并了革新派所有遗产——无论是有形的资源还是无形的声望——且毫无制衡的紫阳宗、丹鼎阁、无极剑派联盟，对天师府而言绝非福音。

那将意味着，天师府这超然中立的地位将受到前所未有的挑战与挤压，甚至可能被卷入更危险的漩涡。

一个僵持的、相互消耗的云霞山战场，才更符合天师府的长远利益——既能持续牵制、削弱陈炎一方的锐气和力量，又能为内域其他观望势力留下反应与抉择的时间。

更何况……

他的女儿，沈瑶，就在那座被围得铁桶一般的山上。尽管父女因理念不合早已疏离多年，当年激烈的争吵与沈瑶的毅然离家仍历历在目，那份血缘牵连与深藏的关切却从未真正断绝。

他得到的情报显示，山上战况惨烈，攻防每日都在见血。

作为一个父亲，他无法想象，若云霞山防线崩溃，沈瑶会遭遇什么。

\vspace{2em}

数日后，天师府总坛传出的命令，便带上了这多重考量混合后的温度。沈天穹以“苍生罹难，修士亦为生灵，不可坐视伤亡惨绝”为由，正式对外宣布，将对云霞山交战区域进行“人道援助”。

这个理由确实冠冕堂皇，任谁也挑不出大的错处，完美地遮掩了其下涌动的战略算计与私心牵挂。

命令下达，天师府这架庞大而古老的机器开始以特有的效率运转。一队队打着天师府杏黄旗号的车队，从总坛及各处分观出发。这些车队主要由未经历练的低阶弟子和谨慎挑选的世俗忠仆组成，车辆普通，护卫力量也仅足以应付寻常匪患，摆明了姿态——我们只是运送救济物资。

车上装载的，是“用以救治伤患、保全元气”的基础物资：品阶不高但数量可观的止血散、养气丹，洁净的灵纱布，以及大量不易腐坏、能勉强补充灵气消耗的低灵谷物。

它们无法决定战局，却足以维系最基础的生命线。

对于兵精粮足、后方补给源源不断的守旧联盟而言，这些“低品”物资不过是隔靴搔痒的锦上添花。他们接收得漫不经心，负责查验的军需官甚至略带嫌弃地检查着丹药的成色，挥挥手便让人将东西丢进庞大的后勤仓库角落，仿佛接收了一样碍眼的累赘。

然而，这些打着杏黄旗的车队试图靠近云霞山时，情形便截然不同了。

守旧联盟的外围哨卡如临大敌，通往山区的每一条小路都已被严密封锁。车队在距离防线尚有数里之地便被凶神恶煞的巡逻队拦下。

“站住！天师府的车队？此地乃军事禁区，闲杂人等不得靠近！”一名身着紫阳宗服饰的小头目按着剑柄，眼神不善地扫视着车队。他身后的士兵们早已结成简易战阵，灵力隐隐流转，封锁了前路。

带队的天师府低级执事是个面相憨厚的中年人，他陪着笑脸，取出盖有沈天穹印鉴的正式文书：“这位道兄，我等奉府主之命，运送人道援助物资，救治伤患，此乃文书凭据，符合天道常伦，还请行个方便……”

“方便？”小头目嗤笑一声，粗暴地打断他，“谁知道你们车里装的是丹药还是爆符？是粮食还是兵刃？前线正在鏖战，谁知道你们是不是打着援助的幌子，给山上的逆贼输送给养？退回去！”

类似的刁难，在各个方向的哨卡反复上演。车队被无故扣留盘查数个时辰，货物被要求全部卸车，一粒粒谷物、一颗颗丹药都被反复用神识探查甚至刮开检验；道路“恰巧”被落石或临时布置的简易法阵阻断；通关文书被挑出各种格式或措辞上的“问题”……

守旧联盟的基层军官们，显然得到了上峰的某种默许，竭尽所能地给这些杏黄旗车队制造麻烦，拖延、阻碍物资的送达。

但他们终究不敢真的撕破脸皮。

当带队执事据理力争，甚至抬出府主沈天穹的名号，并暗示若物资无法送达、伤患死者众，天师府必将质询其违背“天道好生”之责时，那些气焰嚣张的小头目们往往脸色变幻。

他们接到的指令，也确实是“刁难”与“拖延”，“尽量阻止物资上山”。

真正阻挡所有物资，等同于公然打沈天穹和整个天师府的脸。沈天穹本人作为法相境修士的强悍实力，像一座无形的大山，压在这些执行具体任务的修士心头。

他们可以刁难，可以扣留一部分，却绝不敢真的下令攻击车队或彻底没收所有物资。

于是，在反复的拉锯、交涉、妥协之后，总有一部分车队，在经历了重重盘剥与“损耗”后，被允许通过最后一道关卡。

对于被重重围困、资源如沙漏般飞速流逝的云霞山而言，这些历经千难万险才得以流入的援助，其意义远不止于物资本身。

当一袋袋染着尘土甚至血污的低灵谷物，一箱箱品相普通的止血散和养气丹被运回山上时，引发的确是绝境中难以言喻的激动。

久旱逢甘霖，他乡遇故知！

他们并未被外界彻底遗忘和抛弃，有一条极其微弱却真实的生命线，正穿透铁桶般的封锁，连接着山外！

叶牧谦与沈瑶对此心照不宣。沈瑶抚摸着短刀冰凉的刀柄，望向山外遥远的天师府方向，清冷的眸子里情绪复杂难明。

她知道父亲此举，公私掺杂。既有对其理念某种程度上的认可与对平衡局面的算计，也必然少不了对她这个叛逆女儿的担忧。

这份来自冰冷规则与家族隔阂之外的、别扭而隐晦的关切，让她心中某处坚冰，悄然融化了一丝裂缝。

尽管杯水车薪，但这股细微却持续不断的“输血”，实实在在地帮助山上撑过了最初、也是最艰难的消耗阶段。

它让更多重伤员在简陋的救护所有了生还的机会，让筋疲力竭的战斗修士在轮换的间隙，能多咽下一口蕴含灵气的食物，多炼化一颗丹药，从而多恢复一丝宝贵的、可能决定下一次生死交锋的元气。

这是在绝望的钢铁壁垒之中，来自冰冷战争规则外的一丝微弱却真实的暖意，它维系的不仅是生命，更是摇摇欲坠却始终不肯熄灭的希望之火。

\vspace{2em}

过了近一个月。

接连强行冲山的守旧联盟精锐的伤亡、不眠不休推算阵法演进的阵法师的努力，终于使得他们在九宫迷踪的第一重迷雾中，勉强撕开了几条断断续续的通道，总算爬上了半山腰。

然而，这里的迷雾固然变得稀薄，视野稍清，迎接他们的却也不是预想中的开阔地，而是更加赤裸、更加森严的死亡壁垒。

九宫迷踪阵第二重变：借阵法和预先建构的有利地形，迟滞敌人，增强自己。

地形已被改造得面目全非，山路崎岖破碎，迫使他们的百人级大型战阵彻底瓦解，化整为零。

空气中，弥漫着泥石、钢铁与隐隐血腥混合的气息。

就在这里，在这片被改造过的山林里，游击战的铁律开始用最残酷的方式，书写自己的注脚。

叶牧谦的命令非常简洁高效：

——一发子弹消灭一个敌人；

——打一枪换一个地方。

其余的，自由出击。

最初的战斗，并非所有人都能立刻适应。前一句话非常好理解，是告诫他们要节约灵力，应对持久化的战役；后一句话，是什么意思？

对“打一枪换一个地方”这句话不乏反感之人，特别是那些后来加入万仙盟、以勇猛凶悍著称的草莽散修。

一个叫石虎的灵枢境巅峰散修，带着他两个同样脾气火爆的结拜兄弟，被分配在吴刚负责的一片乱石区。

第一次遭遇敌情，是三人的巡哨小组发现了一小队五名紫阳宗斥候，正沿着一条干涸的溪床小心翼翼地上探。

石虎当时就红了眼。“才五个！老子一人就能劈了他们三个！弟兄们，上，拿了这头功！”

他不等队友完全响应，便如猛虎出闸般从掩体后跃出，手中鬼头刀抡圆了，带起一片腥风，直扑最近的那名斥候。

他的突袭确实迅猛，一刀便劈飞了那名斥候匆忙格挡的长剑，随即把他砍倒在地。他的两个兄弟也嗷嗷叫着冲了出去，凭个人勇武各自与剩下几个修士缠斗起来。

战斗起初很顺利，他们凭借个人武力，瞬间压制了对方。

然而，那名被重创的斥候在失去生命前，拼死捏碎了一枚示警符箓。尖锐的灵力啸音划破山林。

“大哥！他们叫了援兵！快走！”一名兄弟格开对手的剑，急声喊道。

“走什么走！还剩两个，宰了他们再走！”

杀红了眼的石虎看着眼前剩下的两个惊恐后退的斥候，觉得煮熟的鸭子不能飞了。

就是这短短几个呼吸的贪功恋战，要了他们的命。

溪床上方，一处看似天然、布满苔藓的岩壁后，突然转出整整一个十二人的紫阳宗快速反应小队，为首的是一名玄脉境开脉期的修士。冰冷的箭矢与灼热的火球瞬间覆盖了溪床区域。

石虎的怒吼戛然而止，他被三支破灵弩箭贯穿了胸膛。他那两个兄弟，一人被火球正面击中，化作焦炭；另一人想逃，却被那灵枢境修士御使的飞剑追上，从后心透入，钉死在溪床的石头上。

他们消灭了三名敌人，却赔上了自己三条命，还暴露了这片区域的防御存在。

消息传回，吴刚脸色铁青。

“我方四十七号哨位已经暴露，短期内不得前往该哨位！”

类似的情况，在最初的几天里并非孤例。

有的小组伏击得手后，觉得敌人慌乱溃逃，有机可乘，多追了几步，结果落入了敌方后续梯队的反伏击圈。

有的小组则固守着某个视野不错的狙击点，连续射杀了几名敌人后，不肯及时转移，很快便被锁定，招来覆盖性的集团法术轰击，连人带工事被炸上了天。

鲜血与瞬间消逝的生命，成为了最冷酷的教官。

叶牧谦说的确实不算数，但一条条活生生的性命，用他们的牺牲反复验证着叶牧谦那两条看似简单的铁律：

——轻易不出动，出动则至少消灭一个敌人；

——出手之后，无论有没有被人发现、有没有叫来援兵，立即换一个地方潜伏。

散修们开始变了。

适应环境，自然也是散修最常见的美德。

他们可能依旧不习惯这种偷偷摸摸、打了就跑的打法，觉得不够痛快，但他们更不想死。

现实的伤亡数字，逼着他们将质疑、将个人勇武的骄傲，死死压在心底。

\vspace{2em}

于是，真正的、高效的游击开始上演。

一个由一名原云隐别院剑修、一名善使符箓的散修和一名猎户出身的神射手组成的三人小组，已经在一处视野极佳、伪装巧妙的树冠狙击点潜伏了整整六个时辰。

他们下方，是一条被落石半掩的废弃小径。

黄昏时分，一队八人的守旧联盟巡逻队沿着小径走来，队形并不算紧密，神色疲惫。树冠上的神射手，将一支刻着破甲符文的铁箭稳稳搭在长弓上，向其中灌注了灵力，呼吸瞬间变得缓慢而深长。

他瞄准的不是为首那个看似指挥者，而是队伍中间那名唯一穿着丹鼎阁服饰、腰间挂着数个药囊的修士。

“咻——”

箭矢离弦的声音微不可闻，下一刻，那名丹鼎阁修士喉咙便被贯穿，哼都未哼一声便扑倒在地。

队伍瞬间大乱。

“敌袭！在树上！”有人指着箭矢来向大喊。

几乎同时，那名忙着画符的修士从另一侧早已看好的位置，扬手打出了三张符箓——两张爆炸符，一张烟雾符，炸伤了几个人，造成了更大的混乱。

“走！”剑修低喝。

三人没有丝毫犹豫，甚至不看战果。

神射手将长弓背好，抓住一根垂下的藤蔓，第一个滑向数十丈下方早就探好的一处岩缝。符修紧随其后。剑修最后一个离开，离开前，他用脚尖触发了一个旁边的机关。

当巡逻队从混乱中恢复，无大碍的五名修士愤怒地朝着原本的树冠位置施展术法、投掷符箓时，那里早已空无一人，只有燃烧的枝叶。

而他们试图追击时，却又触发了疯狂滋长的铁线荆棘，再次被拖延、刺伤。

另一侧。

在泥泞障碍区边缘，另一个小组演绎了另一种配合。

他们故意露出一个破绽，让一支试图绕行的敌方精锐小队发现了他们“仓促”撤退时留下的痕迹。

那支小队果然中计，兴奋地追入一片表面覆盖着落叶、看似坚固的低洼地。

结果，落入了一个连环陷阱。第一个踩上去的修士脚下的“地面”瞬间塌陷，露出一个不算深但布满黏稠淤泥和隐藏水刺的陷坑。

紧随其后的两人急忙想救，却被埋伏在侧翼的陈千羽甩了四根钢针，钢针深深刺入了他们的脚踝，挑断了其脚筋。两人脱力，立即向前栽倒，于是也掉进了泥坑里面。

与此同时，小组中那名远程支援者，冷静地用一枚低阶爆炸符在陷坑旁引爆，巨大的噪音和灵力冲击让陷坑里的人更加慌乱挣扎。

在泥潭中，挣扎反而会使得自己陷进去；只有冷静下来，才能慢慢上浮。

当然他们在泥潭中还受了皮外伤，那救回去八成也会感染——哦，感染在修仙界不算太大的事情，但也能消耗一些后勤。

剩下三人甚至没看清发生了什么事，只能站在坑边，试图拉起三个掉进坑里的家伙。

陈千羽的组员以迅雷不及掩耳之势把这三个家伙踹进泥坑之后，立即远遁，看都未多看那些在泥泞中惨叫的敌人一眼；借助对地形的熟悉，迅速消失在错综复杂的巨石和湿滑的坡地后。

等敌方的支援赶到，只看到六个在泥坑里狼狈不堪、失去战斗力的同袍，其中三个人已经半死不活，敌人早已无影无踪。

——一发子弹消灭一个敌人；

——打一枪换一个地方。

在一次次生死边缘的实践后，这两条律令被山上每一个幸存者用肌肉记忆刻入骨髓。

那些曾经对此不屑一顾的散修，在经历了最初的惨痛损失，又亲眼看到、亲身参与了这种高效而冷酷的杀伤后，心态发生了微妙而彻底的变化。

一次战斗间隙，几个浑身泥污血渍的散修挤在狭窄的掩体里分食着一块硬饼。

其中一个脸上带着新鲜刀疤的汉子，狠狠咬了一口饼，含糊不清地叹道：“娘的，以前觉得叶先生这打法忒不痛快，现在想想……真他娘的高！就这么一点点磨，看着那些紫阳宗的龟孙子明明人多，却像没头苍蝇似的挨揍，真解气！老子这条命，就是听了令，没贪刀，才捡回来的。”

旁边一个年纪稍大的散修，默默磨着自己因为砍人而略有卷刃的剑，幽幽接了一句：“何止是高明……老子活了几十年，打过不少架，从没见过这种打法。让你有力没处使，有火没处发。”

“叶先生……真神人也！”

最后那一声叹息般的“真神人也”，在狭小的空间里回荡，没有豪情万丈，却充满了一种劫后余生、心服口服的感慨。

鲜血与生存，永远是最有说服力的老师。

当那些不信任的、质疑的人都变成了尸体，而严格遵守这些“简单”规则的人却活了下来，并不断给强大敌人制造着伤亡和恐惧时，所有的怀疑都烟消云散，只剩下对制定规则者近乎敬畏的信服。

白言冰、陈千羽、沈瑶、吴东浩、吴刚等核心人物，便如同定海神针，钉在各处要冲，他们自身就是这种战术最坚定的执行者和示范者。

而整条战线，就在这无数个精干如匕首的小组神出鬼没的袭扰下，变成了一架高效而沉默的杀戮机器。

守旧联盟空有数倍兵力，却感觉自己像是在用沉重的铁锤，愤怒而徒劳地击打着无形无影的流沙和铁蒺藜。

他们磅礴的力量被主动分散、被巧妙引导，每一次蓄力的重击，要么打在空处，要么只能换来零星却绝无间断的精准回刺。

看着同袍昨日还鲜活的面孔，今日便成为一具被拖回的冰冷尸体，或者干脆失踪于山林再也寻不回，一股深入骨髓的无力与挫败感，如同这山间弥漫的、无法驱散的湿冷寒气，悄然侵蚀着攻城大军原本高昂的士气。

\vspace{2em}

山下中军大帐内，陈炎接到前线接连传来的受挫战报与日益增长的伤亡数字，脸色阴沉得能滴出水来。 他意识到，单靠中低阶修士组成的常规军队推进，恐怕真要在这座被对手经营得如同铁刺猬般的山上，撞得头破血流，甚至损兵折将、威望扫地。

速胜的耐心被这座顽山消耗殆尽，他眼中寒光一闪，决定不再保留，动用高端力量强行破局，以绝对的实力碾压一切诡计。

“命令周文彬，率本部真符境精锐，自东侧山脊强行御空飞渡，直扑别院核心区域！”陈炎的声音如同金铁交击，冰冷而不容置疑，“告诉他们，不惜代价，给我找出并摧毁其阵法核心节点！”

同时，他转向身旁侍立的随军高层，厉声指示：“让丹鼎阁的高阶药师与紫阳宗的阵法师即刻配合，启动‘焚天’与‘瘴海’预案！我要这半山腰以上，再无一片完整的树叶，再无一处可资利用的地形！”

即便不能直接杀伤那些隐藏极深的老鼠，他也要用覆盖性的烈焰与毒瘴，将山体烧成焦土，将阵基腐蚀瘫痪！

命令下达，山下营地骤然爆发出数道磅礴凶悍的气息，如同沉睡的凶兽苏醒。

以周文彬为首，总计七八名真符境高手周身灵光暴涨，化作数道刺目流光，如逆射的陨星般拔地而起，悍然刺破云雾，直扑山腰以上！

真符境修士已能较长时间御空飞行，足以规避绝大部分恼人的地面陷阱。

在他们身后，数十名玄脉境融贯期的好手也各展手段，或借助法器，或彼此借力，紧紧跟随，意图凭借这石破天惊的一击，从空中强行撕开一条通往胜利的血路。

与此同时，山下数个早已筑好的高大法坛骤然亮起刺目欲盲的光芒，狂暴暴虐的火行灵力与诡异粘稠的墨绿色毒雾开始疯狂汇聚，规模之浩大，灵压之骇人，让远处山巅观看的守卫者们都不由得面色发白，仿佛看到了天灾降临。

然而，中枢阵眼之内，始终如古松般盘坐的叶牧谦，感受到那自天空与山下同时传来的、足以令寻常见独境修士心神摇曳的恐怖灵力波动，缓缓睁开的双眼之中，却掠过一丝“终于来了”的锐利精光。 他似乎，等待的正是对手按捺不住、祭出高端力量的这一刻。

他双手平稳而坚定地虚按在身前——那里是一个由无数道细微光线精准勾勒出的、惟妙惟肖的云霞山脉灵气流动虚影。

磅礴的精神力与他自己独有的、能与周遭天地产生深层共鸣（至于为什么只有他能引起共鸣，叶牧谦并未能深究明白，于是也不去想它了）的精纯元气，如同决堤洪流，奔涌而出，透过身下与整座大山相连的阵法核心脉络，瞬间沟通、唤醒、引导了整座云霞山沉寂却浩瀚的庞大地气。

“地脉如龙，听我号令……起！”一声清喝，虽不高昂，但仿佛蕴含着律令山川的威严。

轰——隆——

九宫迷踪第三重变！

叶牧谦毕竟出身于现代社会，哪里会不知道防空的重要性！

刹那间，整座云霞山主峰，从山基到山巅，似乎被注入了生命，发出一阵低沉雄浑、仿佛源自大地深处的嗡鸣与震颤。

以叶牧谦所在中枢为原点，一层肉眼难以清晰捕捉、却给人以无比坚韧厚重之感的淡金色灵力屏障，如同一个巨大无比、倒扣而下的琉璃巨碗，伴随着嗡鸣声迅速向上蔓延、合拢，最终将山腰以上的核心防御区牢牢笼罩在内！

这屏障中蕴含着叶牧谦对阵道导引、分流、化解的玄妙领悟。

数位真符境长老联手轰出的、足以融化金铁的赤红咆哮火柱，与撕裂空气的凌厉无匹剑气长龙，几乎同时狠狠撞击在淡金色的屏障光幕之上，但预想中的惊天爆炸与屏障破碎并未发生。

那声势骇人的攻击，如同泥牛入海，其大半毁灭性能量被那层流转着奇异韵律的屏障表面巧妙地偏转、卸开、引导。

炽热的火流被引偏斜射向高空无害消散，锋锐的剑气被分化成数百道细流，汇入山体周边呼啸的狂风。余下的威力，则被屏障均匀分散，传导至下方绵延的山体基石，被厚重的大地无声吸纳。

山下法坛蓄势完毕，释放出的、如同赤红浪潮般的烈焰风暴与遮天蔽日的墨绿色毒瘴之云，在触及这层淡金色屏障时，也同样被大幅削弱、稀释、阻隔，未能对屏障后的山体植被与工事造成预期中焦土千里的毁灭效果。

几乎就在淡金色屏障光华大盛、成功抵御住第一波毁灭性打击的同一刹那！

东侧一处陡峭制高点的鹰嘴岩上，沈瑶的身影如鬼魅般悄然浮现。山风猎猎，吹拂着她束起的长发与紧身的玄色劲装，勾勒出凛然如刀的身形。她手中那柄弧度优美的短刀正在微微轻颤，发出几乎听不见却直抵灵魂的细微嗡鸣。

面对数名试图趁乱从侧翼低空掠过、寻找阵法屏障薄弱点或能量节点的敌人，她不退不避，架势倏然展开——

“千幻流光”！

铮——！

刹那间，以沈瑶所处岩石为中心，方圆数十丈内的雾霭、林间疏漏的光斑、甚至屏障流转的淡金微光，仿佛都被强行抽取、熔炼、再造！

数十上百道虚实交织、真假难辨、绚丽夺目却又杀机凛然的璀璨刀光，凭空浮现！

这些刀光如同被赋予了独立的生命与狩猎本能，化身为一条条灵动无比、轨迹完全无法预测的光之游鱼、疾电或幽魂，时而如仲夏夜的流星暴雨，毫无规律地倾泻覆盖一片区域，逼得敌人手忙脚乱地防御；时而又骤然汇聚，凝成一道凝实到极点的深寒刀芒，如同毒蛇吐信，精准刺向因攻击受挫而出现刹那松懈或灵力衔接空白的敌人要害。

真幻交织，虚实互变，令人心神俱疲。

一名紫阳宗玄脉境巅峰的执事，怒喝着挥动火焰盾牌格开了迎面三道看似凶悍的刀光，却冷不防被一道从脚下阴影中悄然钻出、看似虚幻的刀光掠过小腿，护体灵光竟如纸糊般被割开，血光迸现！

沈瑶一人一刀，凭借对地形的极致利用、阵法屏障的掩护以及自身诡谲莫测的领域，竟如一道无形而坚韧的壁垒，生生牵制、迟滞住了数倍于己的敌方高手，令他们无法专心寻找破阵之法，反而不得不分出大量精力应对这无所不在、防不胜防的袭杀。

云霞山，这块被守旧联盟上下视为大军一到便可唾手可得的“骨头”，出乎所有人意料地，在鲜血与火焰的洗礼中，彻底显露出其内里——这根本就是一块坚硬无比、布满倒刺、并且会主动反击的铁桶！

陈炎后续又陆续调集生力军，不惜代价发动数次规模更大的猛攻，前后历时月余，大小血腥接触战数十场，自身伤亡累积已达到一个令中军帐内气氛凝滞的惊人数字，却始终无法突破那最后一道、仿佛拥有生命的核心防线。

第一重变目前已经被全部瓦解，但那毕竟是无人组成、无人参与的；第二重变和第三重变，山上万众一心，此刻根本破不开！

更要命的是，就算他们试图放火烧山，陈千羽的和畅领域一来，草木一夜之间便又长出来了！

这可真真正正是“野火烧不尽，春风吹又生”！

浩浩荡荡的联盟大军，被牢牢拖在了云霞山陡峭的山麓与山腰，进退维谷，骑虎难下。

山下，速战速决、一鼓定乾坤的意图彻底落空，战争无可避免地滑向了他们起初并不愿意但也并不十分介意的消耗泥潭。

陈炎此刻也有些泄劲了。

他的战略也变了：围而不打。

我给你围起来，让你消耗。

你们总不能不吃饭吧？

沈天穹的输血虽然存在，但陈炎冷酷的计算之下，显然不够山上消耗的！

但他显然低估了叶牧谦等人的韧劲：叶牧谦主持下，冷战期间就在山上开凿了水井、开辟了药田、灵田，辟谷丹（粮食）、基础草药均基本自足；天师府又有输血，虽然不多，但蚊子腿也是肉。

要支撑这么多人，勒紧裤腰带一段时间也并非不可。

但能勒紧一时，还能勒紧一世吗？

山上，尽管防线在不少将士以生命为代价的捍卫下暂时稳如磐石，但每个幸存者都能清晰地感受到，那来自山下无穷无尽压力的迫在眉睫。

改变，无论是对于渴望破局的守旧联盟，还是对于必须寻找生路的“问道途”而言，都已迫在眉睫。但至少，死守云霞山的战略目的，已经完成了一部分。

单纯的防御，无论阵法多么精妙，战术多么灵活，终究是在消耗自身本就有限的“血液”。每一块灵石的耗尽，每一瓶丹药的见底，每一处工事的破损，都在削弱这座孤山抵抗的资本。

\vspace{2em}

叶牧谦作为主帅，比任何人都更清醒地认识这一点。

就在山下攻势最为猛烈、山上承受压力仿佛达到顶点的一个深夜，他独自立于别院后方残破的观星台遗址。头顶的夜空，被己方阵法流转的淡金灵光与远方敌军营地连绵的篝火、法坛光芒映照得一片混沌，不见星辰。

然而，前世那些波澜壮阔的记忆碎片，却如同穿透了无尽时空的璀璨星辰，清晰无比地倒映、串联起来——关于在强大敌人围困下依然屹立的红色根据地，关于以弱势兵力周旋并打破一次次“围剿”的惊人战术，关于深入敌人统治腹地、以点点星火最终燎原的传奇。

“坚守要点，以空间换取时间……发动群众……陷敌于人民战争的汪洋大海……”

他低沉的自语声消散在夜风里，轻不可闻，但那双望向山下广袤黑暗地域的眼眸深处，却骤然闪过一丝足以劈开迷雾的决然锐光。

僵局必须打破，而破局的关键力量，不在山上这有限的不到一万人中，恰恰在那山下——那片被守旧联盟视为统治基石，却又在实际上被他们视为可以随意榨取之草芥的、广袤而沉默的底层疆域之中。

\vspace{2em}

翌日，一次参与者极少、保密级别最高的核心会议，在云隐别院深处一座被层层阵法灵光彻底隔绝的静室内召开。空气中弥漫着聚灵阵运转的低微嗡鸣与紧张的气氛。

叶牧谦在一块平整的青石板上，以元气光线勾勒出一幅简陋却关键的内域势力分布草图。他的手指，稳稳地点在了代表着云霞山外围，那一片片连绵的、标注着中小型城镇、散修聚集点和凡人村落符号的区域。

“固守云霞，是为中流砥柱，不可有丝毫动摇。”叶牧谦的声音平稳而有力，仿佛在陈述一个天地至理，“这里将会，并且在未来一段时间内，一直吸引敌军主力。它更是整个‘问道途’的主心骨、精神旗帜，是我们的基本盘，不能放弃。”

他的目光转向侍立一旁的两位大弟子，“白言冰，陈千羽，你二人随我继续坐镇中枢，白言冰统筹防务，临机决断；陈千羽保障救治，维系生机。山在，人在，旗就在。”

白言冰与陈千羽没有丝毫犹豫，同时抱拳，沉声应诺。白言冰眼中是剑锋般的坚定，陈千羽眸子里则是柔韧如青藤的沉静。

然而，问题如同阴云悬在每个人心头。如此被动固守，与敌人拼消耗，哪里能看到胜利的曙光？革新派一方的后勤储备与资源底蕴，远远无法与树大根深的守旧联盟相比，拖延下去，最终的结局几乎是注定的慢性失血而亡。这不仅是所有人的疑问，也正是叶牧谦近期全部思考所聚焦的核心。

“真正的决胜之机，不在山上，”叶牧谦的手指从云霞山图标移开，稳稳地划过草图，落在了那些标志着守旧联盟腹地、但其直属统治力量相对鞭长莫及或较为薄弱的区域，“而在于此处。”

他的目光抬起，落在另外两人身上：“沈瑶，吴刚。”

被点名的两人身躯瞬间绷直。沈瑶清丽的面容上，眸光冰澈凛冽，如同雪原上即将出鞘的剑；吴刚则收敛了平日里那副万事皆可调侃的玩世不恭，粗犷的脸上流露出罕见的郑重与肃穆。

“你们二人，各自挑选一支绝对可靠、精干灵活的小队，人数贵精不贵多。” 叶牧谦的手指关节轻轻敲击着青石板上的那些区域，“目标，不要强行攻打任何重镇要塞，而是潜入这些地方。”

他的指尖重重一顿，仿佛要将决心钉入石板：“那里，有被宗门与家族层层加码的苛捐杂税压得脊梁弯曲、喘不过气的散修；有面朝黄土背朝天，辛勤劳作一年所得大半被盘剥、仅能勉强糊口乃至挣扎在饿死边缘的凡人百姓；还有无数因出身、机缘或无背景，而对现状满怀愤懑却又敢怒不敢言的失意者。”

叶牧谦的语调，带着一种洞悉本质的寒意与热度，“他们是守旧联盟维持统治、汲取养分的根基土壤，但也恰恰可能是这块看似坚固巨石上，最易松动、最先崩裂的一块砖。”

他详细阐述了行动的方略核心：依托对云霞山周边地形地貌的绝对熟悉，借助“九宫迷踪阵”尚未被敌军完全勘破的变幻特性作为掩护，精确计算敌军巡逻队伍的间隙，选择夜色最深浓、戒备最易松懈的时刻，如同水银泻地般悄然渗透出去。

不求一次歼灭多少敌人，不求破坏多少后勤运输，更不贪图占领一城一地。

唯一的核心目的，是将“问道途”追求公道、反抗不公的理念，将那“人人皆可凭自身努力追寻大道、改善生计”的希望火种，尽可能广泛地播撒到那些沉闷而压抑的土地中去。

“具体如何行事，你们可根据实际情况临机决断。”叶牧谦的目光如同实质，扫过沈瑶和吴刚，“但必须牢记我们的根本行动纲领：发动群众，建设组织，积蓄力量。”

他格外强调，“其中，‘发动群众’是重中之重，是源头活水！这场战争注定旷日持久，从长远来看，无数凡人百姓与低阶修士，既是我们的根基所在，也是他们自身改变命运的希望所系！”

“我们能做且必须去做的，”叶牧谦越说越兴奋，语调也越来越激昂，“就是让更多的平民百姓、低阶修士，从心底里理解我们、认可我们、最终愿意支持并追随我们！”

他的视线最终牢牢锁定沈瑶与吴刚，字句如锤：“我要你们，成为插入敌人看似稳固之后方的两把利剑！一把，斩断其汲取民脂民膏、输往前线的血管；另一把，点燃其统治根基之下，那早已遍布干柴的星星之火！”

这绝非一时冲动的冒险，而是叶牧谦结合前世智慧与此世现实，经过深思熟虑后得出的破局之策，得益于他初入内域时，在青石城那长达一年的生活与细致观察。

那一年，让他明确了一个看似简单却至关重要的道理：这个世界虽存在移山倒海的超凡个体，但在大规模战争状态下，高阶修士很大程度上被“兵器化”、脱产（甚至平常情况下也大多是脱产的），而维持这支“兵器”运转的庞杂后勤——从粮秣布匹到基础矿材，从简单工事修筑到短途物资转运——其真正的主力，恰恰是数量庞大的凡人与低阶修士。

也正因洞悉了这一战争维度的本质，他才敢于借鉴并化用那经过历史检验的智慧，坚定地走向这条发动最广大底层力量的群众路线，开辟决定胜负的第二战场。

行动计划迅速在高度保密中敲定。

大多数人可能并不理解为什么修士打架要去管平民百姓死活。但沈瑶有着长达十年的草根经历，吴刚也在万仙盟做过很多低阶散修的工作，所以关于叶牧谦所言“发动群众”，他们是有经验的，也是发自内心认同的。

沈瑶精心挑选了十人，其中多为原无极剑派心剑脉中剑法以轻灵迅捷见长、尤其擅长隐匿潜伏的弟子，再搭配数名在加入万仙盟前便是独行游侠、极擅野外生存、陷阱布置与情报侦查的散修老手。

吴刚的队伍则更侧重综合能力与适应性，既有万仙盟中熟悉各地黑市门路、人情往来与社会三教九流的“路路通”，也有像陈默这样早期转修问道途、能作为鲜活榜样吸引同类的修士，还包括两名思维敏捷、口齿伶俐、善于在不同场合进行说服与鼓动的特殊人才。

月黑风高，正是潜龙腾渊时。

在云霞山防御部队又一次浴血击退敌军组织的夜袭、战场迎来短暂死寂的间隙，两条比夜色更加深邃、更加沉默的“细流”，借助山间尚未散尽的雾气与复杂崎岖地形的完美掩护，从守旧联盟包围圈两个因地形限制而形成的、极为细微且不易察觉的结合部缝隙中，悄然滑了出来。

沈瑶一袭不起眼的深蓝近黑劲装，将姣好的身形完全融入阴影。她甚至未曾全力催动见独领域，仅仅让其维持在周身尺许，完美地扭曲光线、吸纳气息。

她如同夜幕中飘忽的领路幽魂，身后两支小队成员个个屏息凝神，将一切行动声响降至最低，如同真正的鬼魅般紧跟着她，穿行在林木与陡崖的阴影之间。很快，那旌旗如林、灵光冲天的连绵营寨被远远甩在身后，模糊成一片背景光晕。

他们的面前，是沉沉睡去的广袤山林，以及山林之外，守旧联盟统治下那些光影交织、苦难与麻木并存的灰色地带。

两队人马在一个早已标记好的三岔口悄无声息地分开，互递一个坚定的眼神后，便向着不同方向，隐没于更深沉的黑暗之中。

\vspace{2em}

沈瑶带领的小队，彻底化为了一柄经过千锤百炼、敛尽所有华光的漆黑匕首，沉默、精准、致命地刺向敌人后方。他们昼伏夜出，完全避开官道与任何稍具规模的城池，专挑那些小路行进。

目标极其明确：那些由守旧联盟附属中小宗门或地方家族控制、驻防力量有限、看似不起眼，却实际承担着为前线大军输送物资、同时层层盘剥底层以供自身享乐的城镇节点。

七日后，经过周密迂回，他们抵达了第一个目标——位于紫阳宗势力范围边缘地带的灰谷镇。此地因出产一种用于低阶法器淬火增韧的辅料“灰纹铁”而得名，镇民结构简单，多为在矿洞中讨生活的矿工、往来运输的贩夫走卒，以及少数依附这条产业链生存、修为低微的散修。

镇守此地的，是紫阳宗某外门长老的一个姻亲家族，修为最高者不过玄脉境中期，但凭借后台，在镇内作威作福，苛捐杂税名目繁多，征敛无度，是标准过头的土皇帝。

沈瑶没有丝毫急躁。她亲自带领两名最擅长潜行侦查的队员，如同一缕青烟般在镇内外潜伏、观察了整整三天，摸清了守卫懒散的巡逻规律、镇内几处关键仓库与货栈的具体位置与守备情况，以及镇上矿工、散修们主要聚集闲聊的几个茶摊、酒肆。

她清晰地看到，因为战争，紫阳宗对灰纹铁的收购价格被压得更低，而经由联盟商会控制的粮食、盐、粗布等生活必需品的价格却飞涨。镇民们脸上菜色更重，眼神麻木中压抑着怒火，私下怨声载道，但面对镇守家族那几个嚣张跋扈的护卫和背后的紫阳宗，无人敢带头反抗。

行动日，选定在镇上一次旬日小集的日子。正午时分，狭长的集市街道略显拥挤，充斥着讨价还价声、牲畜叫唤和粗重的咳嗽声。沈瑶命令所有队员——也包括她自己在内——换上与当地矿工、农夫无异的粗布短打，甚至故意蹭上些尘土泥污。

大部分队员化整为零，悄然散布在集市外围关键路口，如鹰隼般警戒。她自己则仅带两名最得力的队员，步履自然地穿过人群，径直走向位于集市中段、那间由镇守家族开设、兼做着收税与强买强卖勾当的最大粮铺。

粮铺掌柜是个脑满肠肥的中年人，正尖声呵斥着一个蜷缩在柜台前、颤抖着掏不出几块下品灵石购买全家半月口粮的老矿工，唾沫星子几乎喷到老人脸上。忽然，掌柜浑身肥肉一僵，感到脖颈侧面传来一阵深入骨髓的冰凉，一柄弧度优美、却散发着森然死意的幽暗短刀，已无声无息地贴在了他的皮肤上。

“所有人，不许动，原地蹲下！”

沈瑶清冽的声音并不高昂，却奇异地压过了集市所有的嘈杂，清晰地钻入每个人耳中。与此同时，她身上那股修士特有的、凝练如精钢的气息微微释放了一瞬。

如同寒风掠过，顿时让粮铺内几名仅有灵枢境修为、平日里狗仗人势的护卫双腿发软，脸色煞白，僵在原地动弹不得。

集市上的人群先是一静，随即爆发出本能的骚动与惊呼，但在外围队员冰冷而有序的维持手势，以及沈瑶身上那似有若无、却足以让普通人灵魂战栗的凛冽气场笼罩下，骚动迅速被压制下去，化为一种死寂般的、充满恐惧与惊疑的安静。

沈瑶看也不看那抖若筛糠、面如土色的掌柜，她的目光平静地扫过集市上那一张张或麻木、或惶恐、或畏缩、又或在不经意间流露出一丝好奇与期盼的脸庞——那是被生活重压碾磨过的底层百姓最真实的面容。

她深吸一口气，朗声开口，声音在元气的包裹下传遍半个集市：

“我乃云隐别院，叶牧谦叶先生座下弟子，沈瑶。今日来此，不为杀戮，不为劫掠，只为替各位父老乡亲，取回本就属于你们的活命之物！”

话音未落，她手腕看似随意地一抖，一道凝练的元气刀光掠过，粮铺柜台上那本记录着巧取豪夺与高利盘剥的厚重账册，应声被挑飞至半空，纸张哗啦散开，如同片片控诉的雪花。

几乎在同一时刻，她身旁一名队员如同蓄势已久的猎豹般窜出，身法快得只在众人眼中留下一道残影，几道精准如尺量般的剑气“叮叮”数声轻响，粮铺通往后院仓库那扇厚重木门上的精铁大锁应声断裂。

第三名队员则早已如同鬼魅般贴近了闻讯从隔壁酒楼仓皇冲出的几名镇守家族护卫，简单利落的几记手法，便让他们瘫软倒地，暂时失去了行动能力。

“去，打开仓库！将里面囤积克扣的粮食、盐、布匹，全部搬出来，就在这集市空地上，按户按人头，公平分发！”沈瑶的命令清晰果断。

她的队员们毫不犹豫地执行。更令人动容的是，那几个原本被掌柜呵斥、蹲在墙角的老矿工，此刻眼中骤然爆发出一种近乎凶狠的光芒，那是对生存最本能的渴望被点燃的火焰！他们低吼一声，仿佛瞬间年轻了十岁，跟着问道途的队员们一起，冲进了那座平日里他们连靠近都不敢的仓库。

当第一袋沉甸甸的、散发着谷物清香的麻袋被扛出来，随后是第二袋、第三袋……当一匹匹虽然粗糙却厚实的土布、一罐罐洁白如雪的盐巴被陆续搬出，在集市空地上逐渐堆积成一座小山时，死寂的人群中终于爆发出了难以置信的、混杂着狂喜与哽咽的低呼与骚动。

许多人的眼睛死死盯住那些物资，眼眶迅速泛红。

沈瑶知道，火候已到。她的声音再次响起，依旧清冽，却注入了一种铿锵的、直指人心的力量：

“乡亲们看好了！这些粮食、这些布匹、这些盐，哪一样不是你们的汗水换来？哪一样不是守旧联盟那些高高在上的老爷们，从你们嘴里夺去、从你们身上剥下，用以填满他们自己的仓库，供养前线那些用来继续压迫更多像你们一样百姓的军队？！”

她的话语如同匕首，剖开了血淋淋的现实，“他们何曾将你们当人看待？在他們眼中，尔等不过是可以随意驱使、榨取、丢弃的牛马与蝼蚁！”

她停顿下来，目光如电，扫过人群。她看到，许多双原本麻木的眼睛里，开始燃起压抑已久的火焰；许多紧抿的嘴唇在颤抖；一些年轻人捏紧了拳头，指节发白。那是积郁多年的悲苦与愤懑，被这番话说穿、点燃后的共鸣。

“我师叶牧谦，观天地不公，悯众生之苦，另辟‘问道途’！”沈瑶的声音陡然拔高，带着一种宣告般的穿透力，“此道，不假外求，不拜神佛，只修己身！无需依赖大宗门垄断的昂贵灵石、稀缺丹药、天价法器，凡人亦可凭自身毅力，引气入体，踏上修行之路！即便是资质寻常的低阶修士，修习问道途后，其劳作之能、体魄之健、乃至心志之坚，都远非寻常灵枢径修士可比！这意味着什么？”

她自问自答，话语斩钉截铁：“这意味着，跟着我们，你们或许无法立刻飞天遁地，但你们能凭自己的双手，开垦更多荒地，收获更多粮食；能凭自己的力气，从事更繁重的劳作，换取更安稳的生活；能让你们的子女，有机会不花一个冤枉钱，就接触到真正的修行法门，拥有改变命运的可能！而不是世世代代，被锁死在矿洞、田垄和最低贱的活计上，永无出头之日！”

这番话语，没有任何虚无缥缈的大道理，而是用最朴素、最直接的方式，描绘了一个触手可及的、关于“更好生活”的图景。

对于这些挣扎在生存线上、对高阶修士世界如同仰望星空的底层百姓而言，什么长生久视、什么大道法则都太过遥远；唯有实实在在的粮食、衣物、盐巴，以及让下一代人过得比自己稍好一点的希望，才能真正撼动他们坚固如铁石的心防。

“守旧联盟为何要不惜代价，围困云霞山，打压我等？”沈瑶的语调转为凌厉的质问，“正是因为他们恐惧！恐惧一旦人人皆可凭自身努力修行、改善生活，无需再仰仗他们的鼻息、购买他们垄断的资粮，他们那作威作福、骑在所有人头上吸血的特权，就将如同烈日下的冰雪，彻底崩塌消散！”

“我们一退再退，他们却步步紧逼，甚至不惜设下埋伏，屠戮我同道！试问，我‘问道途’求一个公道，何罪之有？你们这些终日劳作不得温饱的百姓，又何罪之有，要承受这般无休止的盘剥与轻贱？！”

此时，物资分发已经在队员和一些胆大起来的镇民主持下，开始有序进行。

当第一个头发花白的老妇人用颤抖的双手接过一小袋粮食和一小块盐巴时，浑浊的泪水瞬间滚落；当一个面黄肌瘦的孩子紧紧抱着一块粗布，脸上露出罕见的、属于孩童的笑容时……整个集市的气氛变了。

先前对沈瑶这位“修士”的恐惧与隔阂，在她身穿粗布衣、口吐百姓言、行事为百姓夺回生存之资的对比下，迅速转化为一种难以言喻的信任与激动。

那是一种最质朴的认知：这个“仙子”和以前见过的、那些高高在上、冷漠索取的人们，不一样。

沈瑶敏锐地捕捉到了这种变化。她向人群中几个一直紧紧盯着她、眼中除了激动还有某种更深层渴望的年轻散修招了招手。

这几人修为低微，多在灵枢初、中期徘徊，衣衫陈旧，显然是镇中最底层的那批散修。她接过队员递来的几本纸张粗糙、装订简易的薄册子。

“此乃我师亲撰、简化优化的问道途《养气篇》最基础入门法诀。”沈瑶的声音缓和了些，带着一种郑重的意味，“文字浅白，路径清晰，于资源要求降至极低，唯重心性与毅力。此乃正道之基，虽简，却直指修身自强之本意。有意者，可自取观之、习之。能否入门，踏上此途，全在个人抉择与坚持。”

她将册子亲手递给那几名眼神灼热的年轻散修，又额外留下几份，放在分发物资的木台上。

“我辈修士，生而在世，当有自强不息、不仰人鼻息之志！而非浑浑噩噩，任人鱼肉宰割！”

一位衣着最破旧、手上满是矿洞污渍却洗得发白的年轻散修，几乎是颤抖着接过册子，急不可耐地翻开第一页。扉页之上，并非什么高深功法总纲，只有十个大字：

“仰不愧于天，俯不怍于人。”

年轻散修浑身一震，呆呆地看着这十个字，眼眶瞬间通红。他身边另一个同伴凑过来看，也是身躯微颤，低声喃喃：“不愧本心……当今之世，能做到的，有几人？”

修士与凡人百姓的关注点，在此刻微妙地分化开来。

大多数百姓接过救命的粮食布匹，感激涕零，心中记住的是“叶先生”、“沈仙子”给了他们活路，对那十个字或许感触不深，只觉得是“好人该有的样子”。

而这些挣扎在修行底层、看多了蝇营狗苟、自身也常因资源匮乏而不得不做出些违心之举的低阶散修们，对这十个字的共鸣却如同惊雷，直击灵魂深处。

它触碰的，是他们身为修士，却活得不如意、甚至时常自我质疑的那份尊严与骄傲。

整个分发过程高效而迅速，持续了不到一个时辰。

在镇守家族可能的主力修士闻讯从烟花之地或修炼静室匆忙赶来之前，沈瑶小队早已带着缴获的部分关键账册——作为守旧联盟盘剥的罪证——和少量便于携带的灵石，如同他们出现时一样，悄然撤离，瞬息间便融入镇外山林，消失得无影无踪。

留给灰谷镇的，不仅是家家户户或多或少分到的、足以支撑一段时日的活命物资；不仅是那几本在暗地里被无数人渴望、偷偷传抄临摹的《养气篇》入门册子；更重要的，是“沈瑶”“叶牧谦”“问道途”“仰不愧于天，俯不怍于人”这些名字、话语和理念，如同巨石终于投入了这潭被压抑已久的死水，激荡起层层叠叠、再难平息的思想涟漪与渴望的波澜。

关于云霞山的故事，关于另一种活法的可能，开始在这个小镇的街头巷尾、矿洞内外，以各种被添油加醋却核心不变的版本，秘密而顽强地流传开来。

此后月余时间里，沈瑶带领这支精干小队，如同技艺高超的弈者，在守旧联盟广袤的统治棋盘上，落下了一枚枚看似微小、却位置关键的棋子。

他们如同飘忽不定、却总能精准命中要害的幽灵，连续“光顾”了四五个与灰谷镇情况类似的城镇据点，有时会选择在当地守旧势力头目正在大摆筵席、醉生梦死时突然现身，以雷霆手段控制全场，在众目睽睽之下揭露其奢靡与底层困苦的对比；有时则会暗中联系当地一些早已对压迫心怀不满的小型散修团体或受排挤的家族旁支，里应外合，事半功倍；在一次特别行动中，他们甚至成功伏击了一支押送队伍，解救了一群被守旧联盟强行征发、准备送往最危险的前线地段充当肉盾“炮灰”的散修苦力，并将他们妥善安置、吸纳……

随着这些行动如水波般扩散，“问道途”以及其核心人物叶牧谦的形象，在广大的底层区域，开始以一种复杂而多元的方式被认知、被谈论。

对于绝大多数挣扎在温饱线上的凡人百姓而言，“叶先生”“云隐别院”这些词汇，最初或许模糊而遥远。但经由沈瑶等人一次次实实在在分发粮食、布匹、盐巴的行动，“问道途”逐渐与“能让我们吃上饭”、“不欺负我们”、“替我们出头”这些最朴素、最根本的利益诉求画上了等号。

他们或许不懂“不假外求”的深奥，也难以理解“仰不愧天”的境界，但他们心中有一杆最朴素的秤：谁对他们好，他们就记着谁的好。叶牧谦的形象，在他们口中传颂时，常常被赋予一些朴素的想象，有时是“下凡救苦的神仙”，有时是“法力无边、专打坏蛋的大侠”，有时甚至就是“一个特别厉害、心肠好的读书人”。

这种认知虽然不够准确，甚至有些神化或简单化，却无比真实地反映了底层民众最直接的感激与期望。

所有平民百姓，实际上都渴望一个能带来公平、让他们活下去的救星或青天。

而对于数量稍小、处境稍好但也极其有限的低阶散修，以及部分出身寒微、在守旧联盟体系内不得志的低阶修士而言，他们对叶牧谦和“问道途”的看法则要复杂、深刻得多。沈瑶留下的那本《养气篇》册子扉页上的十个字，如同黑暗中的一盏孤灯，照亮了许多人内心不曾明言的委屈与坚守。

“问道途”宣称的“无需依赖垄断资源”“仅凭自身努力亦可修行进阶”，更是切中了他们修行路上最大的痛点——资源匮乏与上升无门。

叶牧谦的形象，在他们心中更接近于一个“开宗立派、为吾辈寒微者开辟新路的大贤先师”。

他不仅提供了一种可能更公平、门槛更低的修行路径，更重要的是，他提供了一种精神上的认同与尊严。他们讨论叶牧谦时，会探讨他那奇特的“养气”“见独”境界究竟有何奥妙，会争论“问道途”与“灵枢径”在根本理念上的差异，会羡慕白言冰、沈瑶等人能追随这样的师长。

叶牧谦之于他们，既是现实摆脱困境的希望，也是一种理想人格的投射。

\vspace{2em}

吴刚那面采取的行动，和沈瑶则大相径庭。

如果说沈瑶行动在明，如同撕裂黑夜的雷霆与烈火，以直接而凌厉的方式发动群众，凭借夺回物资、公开宣讲，在底层百姓心中迅速点燃对守旧联盟的憎恶与对“问道途”最朴素的认同；那么，在另一条更为隐蔽的战线上，吴刚所执行的，则是战略的另一个精妙侧面：团结一切可以团结的力量。

这显然是基于两人截然不同的经历与资源。

沈瑶十年游历，深谙草根疾苦与民间生态；而吴刚的优势，则根植于万仙盟那张早在冷战时期就已存在、并在之前经济封锁中历经考验而变得越发坚韧细密的地下网络，以及极其善于揣摩人心、权衡利弊、在复杂利益格局中找到平衡点与共鸣处的特殊才能。

虽然说先前那个网络在守旧联盟巨象踏步一般的前期推进之中，不少敢于露头的地下网络联络员都被抓到，面临监禁，甚至当场牺牲；但依然有着少部分原地下网络联络员依然保存了自己的上下游联络方式、情报和性命。

不过已经建立过一次地下网络，在原网络的废墟之上重建一个网络，那难度相形之下会稍低些。

吴刚选择的切入点，自然也是散修比例更高、商业活动更活跃、各方势力交错、情况也更为微妙的区域——情况越混乱，乡土气息越浓厚，那么守旧联盟的触角就更难管辖到哪些地方，重建网络、团结地方的操作自然也相对容易一些。

他率领小队潜出重围后，抵达的第一个重要节点，便是这样一个地方——名为“黑石集”的大型散修贸易据点。

\vspace{2em}

黑石集坐落在一片盛产低阶炼器矿石与药材的山脉出口。建筑杂乱无章，街道上弥漫着矿石粉尘、药材腥气与汗臭混合的味道。

这里名义上受一个依附于丹鼎阁的中等势力管辖，但天高皇帝远；真正的秩序维持与利益分配，很大程度上依赖于散修们自发形成的潜规则，以及几个在此地盘踞多年、彼此制衡的本土小帮派。这里是信息的集散地，也是资源的流转池，更是无数散修讨生活、碰运气的江湖。

吴刚如同一条回到熟悉水域的游鱼，通过队伍里的本地散修，只用了不到两天，便在一家充斥着劣质酒气与粗豪喧哗的嘈杂酒馆后间，与黑石集本地一位颇有威望、人称“老黑石”的散修头目接上了头。

老黑石是个年长的矿工头目，人如其名，皮肤黝黑粗糙如历经风霜的岩块，一双眼睛却锐利依旧。他打量着一身寻常散修打扮、笑容里带着三分市侩七分诚恳的吴刚，以及吴刚身后那位气质沉静、与周遭环境隐隐格格不入的陈默，并未立刻表态。

吴刚也不废话，寒暄两句后，直接伸手入怀，取出的是一卷看似普通的账册抄录。

他将抄录在油腻的木桌上摊开，手指点向上面几行迥异的数字。

“老黑石，您是明白人。咱们不讲虚的，看这个。”

吴刚的声音压得很低，却字字清晰，“您看仔细了。这一行，是丹鼎阁今年在咱黑石集收购黑须草的‘官定’收购价，明码标价，童叟无欺，对吧？”

黑须草算是一种不算多难培植的灵草，能够代替更常见的凝气草，在黑石集附近地区繁盛。这是黑石集附近部分从事农业散修的主要生产对象。

老黑石虽然是矿工，但是家里不乏种植黑须草的，所以行情也大概了解。

这个价格，和他的听闻差不多。

吴刚指尖一滑，落到下一行：“您再瞧这一行。这是同一批黑须草，被他们运出去，简单炮制一番后，转手卖给前线各宗，或者是炼制成最低阶的回气散后，在联盟控制的大城药铺里的售价。”

老黑石鼻翼微微翕动，盯着那串数字的目光仿佛要将其烧穿。

他仔细地看了一遍账本，却并未发现什么作假之处。

吴刚啧啧两声，摇着头，脸上写满了替人不值的唏嘘，“这还仅仅是一项，要是加上矿业、养殖业，中间的差价……嘿，我粗算了一下，刨去成本，剩下的利润，够把咱们黑石集老老少少、所有散修兄弟的吃喝用度，舒舒服服提升好几个档次，甚至能供养几个有天赋的后生安稳修炼到灵枢后期了。”

老黑石脸色本就黝黑，此刻更显得阴沉了几分。

这些事，他何尝不知？

黑石集的散修们起早贪黑种植的黑须草，冒着被妖兽袭击、矿洞坍塌的风险采集来的药材矿石，大半利润都被上游的丹鼎阁及其附庸家族轻易攫取。他们就像守在泉眼边的挑水工，拼死拼活挑上来二十桶水，自己能留下的，不过区区数桶。

吴刚观察着他的神色，又适时递上一张新的纸，上面是紫阳宗最新下达的征税令。“还有这个，紫阳宗的‘特别军税’，这个月额度比上个月又硬生生抬高三成！理由是‘云霞山战事吃紧，需加大投入’。”

他语带嘲讽，“可咱们得了什么好处？前线死的散修兄弟，抚恤金可有一枚灵石足额发到他们家人手里？咱们黑石集被征调去的后生，是编入了紫阳宗、丹鼎阁的核心战阵享福，还是被充作冲锋在前、试探陷阱、消耗对方箭矢的仆从军、敢死队？”

老黑石腮帮子咬得咯吱作响，眼中闪过一丝痛楚与无力。

这些都是血淋淋的现实，他手下就有几个兄弟一去不回，家人只等到一张轻飘飘的“阵亡通知”和几块打发叫花子的灵石。

更让他内心备受煎熬的是，面对丹鼎阁和紫阳宗的层层压力，他这个本地头目有时也不得不协助征调、催缴税款，明知是火坑，也得推着乡亲们往里跳，背地里不知挨了多少骂，听了多少戳脊梁骨的话。

虽然市侩，但他至少还有良心。

吴刚身体前倾，声音压得更低，却如淬毒的细针，直刺老黑石的心防：“老哥，咱们敞开了说。在紫阳宗、丹鼎阁那些高高在上的大人物眼里，咱们散修是什么？”

随即自问自答，每个字都透着冰冷的寒意，“是耗材！是用完即弃的灵符箭矢！是维持他们战阵运转、可以随意填补进去的数字！是需要时肆意强征暴敛、无用时便弃如敝屣的蝼蚁！”

“可我万仙盟，为何最终铁了心，跟云隐别院、跟叶先生站在一起？”吴刚话锋一转，语气带上了一种同道相认的诚挚，“就是因为看透了这一点！咱们散修，不想永远当耗材，当蝼蚁！”

他指向身旁一直沉默不语的陈默：“陈默道友，您或许听说过。原先也就是个跟咱们差不多的普通散修，机缘巧合得了叶先生指点，转修‘问道途’。您看看他现在，修为精进多少暂且不说，这精气神，这眼神里的笃定，您感受感受！在云隐别院，凭本事说话，有功则赏，叶先生传道解惑，何曾收过弟子一枚灵石的孝敬？”

陈默适时地微微颔首，周身那股沉静而内蕴锋芒的气度，与他朴素的衣着形成奇特的和谐，本身就是最具说服力的无声证言。

老黑石的目光在陈默脸上停留片刻，又回到吴刚身上，眼中的警惕与犹豫终于开始松动，取而代之的是一种强烈的、对现状不甘的共鸣，以及对另一种可能的好奇。

吴刚知道火候已到，不再停留于控诉，而是抛出了实实在在的合作方案：“老哥，不瞒您说，云霞山是被重兵围着，但您也是老江湖，该看得出，那不是一时半会儿能啃下来的硬骨头。 长久拖下去，守旧联盟这般涸泽而渔，四处树敌，内里能支撑多久？反观问道途这边，得道多助，路子越走越宽。 我们不求黑石集的兄弟们立刻明火执仗地造反，那是把大家往绝路上逼。但有些事，咱们可以暗中携手，互惠互利。”

他条理清晰地列出几条建议：黑石集的散修可以暗中将一部分采集到的药材、矿石，以远比丹鼎阁“官价”公平的价格，出售给吴刚联系的秘密渠道；对于联盟下达的征调令，可以采取阳奉阴违、拖延敷衍的态度；在不暴露自身核心的情况下，可以传递一些关于本地守军动向、物资运输路线等非绝密但对问道途了解后方态势有益的情报；甚至，可以暗中关照、保护那些在黑石集内，因听了传闻而对“问道途”产生兴趣、偷偷尝试修炼那本《养气篇》基础法诀的年轻人。

“作为回报，”吴刚的承诺同样具体而务实，“我们万仙盟的网络，可以帮黑石集的兄弟们，开辟一些……嗯，守旧联盟控制区之外的‘贸易选择’。你们需要但被联盟卡脖子的特定药材种子、某些特殊的工具法器图谱、甚至是一些外界的信息，我们可以想办法。 我们还可以分享一些，如何‘合理’规避某些过于严苛征役的小窍门；万不得已时，甚至能为一些兄弟的重要家眷，提供远离这战区边缘的临时安置渠道。”

当然，吴刚深谙驾驭之道，明白纯粹的利诱不足以长久，必须辅之以必要的威慑。

“当然，丑话说在前头。合作归合作，底线不能破。” 他脸上的市侩笑容收敛，目光变得锐利，“若是黑石集的哪位，借着咱们合作的机会，转头去欺压更底层的矿工、盘剥普通百姓，那我的人也会一笔一笔记得清清楚楚。将来若是时移世易，大局有变，我军真的得了势，到了清算的时候，这些人会是什么下场，我可不敢保证。”

最后，他总结道，语气恢复了那种令人容易接受的、权衡利弊的商人式诚恳：“风险可控，利益可见。 他们视咱们为蝼蚁，随时可以牺牲；而我们，却愿视黑石集的诸位为潜在的同道，为将来大业的基石。今日在此埋下合作的种子，浇水施肥，来日若得天时，未必不能生长为庇护一方、让兄弟们活得更有尊严的参天大树。即便最坏的情况，咱们没能成事，对黑石集而言，多一条隐秘的退路和财路，总比只在守旧联盟这一棵树上吊死、被一点点榨干骨髓来得强，对吧？”

这番立足于冰冷利益计算、却又巧妙描绘出远景与尊严的剖析，没有空泛的口号，每一句都敲在老黑石和类似头目们最关切、最疼痛的点上。它揭示了被剥削的现状，提供了切实的替代方案，指明了潜在的风险与远期的希望。

老黑石沉默了许久，最终，重重地将杯中残酒一饮而尽，酒杯落在桌上发出一声闷响。

这，便是应允的信号。

合作，就此在不见光的暗处生根。

虽然这种操作在沈瑶这种磊落的人看来并不光彩，但确实是非常务实而有效的了。

吴刚行事极为周密。

会谈之后，队伍中那名原本就出身黑石集附近、对当地人情脉络了如指掌的散修，便以“游子厌倦漂泊，归乡定居”为由留了下来。

明面上是落叶归根，暗地里，则承担起联络、协调，以及某种程度上的“观察”职责，确保合作沿着既定轨道进行，防止有人首鼠两端或趁机作恶。

另一方面，那个散修钉进黑石集之后，新的情报网络的最初节点自然也形成了。先前因躲避危机而停用、休眠的旧网络，也自然会随着新网络的启用而慢慢复苏。

这张无形的网，开始成为散修们共享信息、协调行动、彼此声援的地下生命线。联盟哪里又加了新税？哪支巡逻队换了凶狠的头目？哪个家族的仓库最近看守特别松懈？

这些零碎却至关重要的信息，开始在网络中悄然流动。

\vspace{2em}

吴刚更深的谋划，在于人的塑造。

他的眼光毒辣如鹰隼，目标不是那些热血上头、叫嚷最响的年轻人，也不是那些早已麻木、逆来顺受的老油子，而是那些本身在一定范围内具有威望、头脑相对清醒、对现状积累了足够不满、内心深处依然渴望改变的“中间层”。

“老刀子”秦厉，便是这样一个被吴刚“相中”的人。

他是黑石集南区矿洞的一个矿工，人如其名，脾气又硬又直。

但老刀子近来越发沉默，眼神里时常流露出一种深沉的疲惫与郁结。他侄子三个月前被紫阳宗的征调令强行带走，说是补充前线辅兵。一个月后，只传回一块破碎的身份玉牌和五块下品灵石的抚恤。他偷偷托人去打听，隐约得知侄子是被派去执行一次危险的试探性冲锋，死得毫无价值，尸骨都没找回。

本来骂了两句问道途，这事情就应该结束了；但他之后又听说，同期被征调的、几个依附于丹鼎阁的宗门子弟，经上下打点，却大多被安排在相对安全的后勤位置。

这件事像一根毒刺，深深扎进老刀子心里。

自此，他开始平等地怨恨问道途和守旧联盟，但更多的是无力。

他知道自己这点力量，对抗不了庞然大物，曾经的热血和硬气，在冰冷的现实面前，似乎只剩下每日收工后，在烈酒中寻求片刻麻木。

吴刚注意到老刀子，是通过那名留在黑石集的联络员。几次“偶遇”和酒馆里的观察后，吴刚判断，这是一个有底线、有担当、且内心痛苦正在寻求出口的人。他需要的不是煽动，而是一个支点，一点火星，以及一条若有若无、却切实存在的路。

第一次接触，不是在酒馆后间，而是在老刀子收工后，独自坐在矿洞外一块大石上，对着夕阳发呆的时候。吴刚扮作一个路过歇脚、收购零散矿石的行商，很自然地坐到了不远处，递过去一袋水，随口聊起了今年矿石的品相难寻，物价飞涨。

老刀子起初戒备，但吴刚谈论的都是最实际的行情，语气里没有高高在上的修士架子，反而像是个精明的生意人，偶尔还骂两句丹鼎阁收购官压价太狠，盘剥得太凶。这些抱怨说到了老刀子心坎里，两人的话匣子慢慢打开。

吴刚没有一上来就提“问道途”——他也知道老刀子对问道途的印象好不到哪去——他只是像一个见识广博的旅人，讲述着沿途见闻：“往南走，过了青岚江，那边的散修日子也不好过，但听说有些地方，散修自己抱团，搞了些小型的‘互助会’，一起接活儿，一起跟收购商讨价还价，虽然也艰难，但听说被盘剥得没咱们这儿这么狠……唉，也是听说，不知真假。”

老刀子听着，浑浊的眼睛里闪过一丝微光，闷声道：“抱团？谈何容易。没个领头扛事的，没点底气，一盘散沙，一吓就散。”

吴刚深以为然地点点头，状似无意地提起：“也是。不过我还听说，云霞山那边的散修，好像有点不一样。那边被围得铁桶似的，但里头的散修听说反而挺硬气。”

“别他妈的跟我提云霞山那边，闹心。”老刀子道。

“哈哈，不是夸他们。我又听说，有个叫什么陈默的，以前也就是个普通散修，现在可是云隐别院叶先生眼前的红人，又没打点又没干啥的，但修为精进不说，管着不少事儿呢。啧，这人跟人，命不同啊。”

“陈默？”老刀子似乎对这个名字有点印象。

黑石集消息杂乱，关于云霞山的只言片语里，好像确实提过这么个人物。他以前只当是谣传，此刻从一个看起来实在的“行商”嘴里说出来，分量似乎重了一点点。

“他不是吹出来的？”老刀子还是不相信。

“谁知道呢，道听途说。”吴刚摆摆手，结束了这个话题。

转而，说起一种能缓解矿工长期劳作筋骨酸痛的廉价药油配方，说是南方某个小地方流传的土方子，材料便宜，就是配伍有点讲究。

他很简单地说了说主要材料和大概比例，然后就像闲扯完一样，拍拍屁股起身，说还要赶路去下一个集子。

老刀子将信将疑，但矿工们的劳损病痛是实实在在的，他留了心。

几天后，他试着按那模糊的方子配了点药油，给几个老伤痛的兄弟用了用，反馈居然不错，虽然比不上真正的丹药，却也能缓解不少痛苦。

这让他对那个神秘“行商”的话，又信了一分。

当然，纵然如此，想要改变一个人对问道途的看法，道阻且长。

这只是开始。

这个行商倒是见多识广，有时候三天来一次，有时候十天半月不来一次；但每次和他闲聊，总能套到些不错的消息。

有时是一批市面紧俏、能有效防护矿尘的低阶面罩材料来源信息；

有时是关于紫阳宗某个税吏喜好和弱点的几句提醒，让老刀子能在交涉时少吃亏；

有时则是一些外界流传的、关于守旧联盟内部倾轧、后勤吃紧的模糊消息；

吴刚有时候也会装作不经意间提到问道途，讲的倒是半真半假，也算客观；又说有个女侠据说出自问道途，这几天四处乱窜，到处破坏敌人统治，乃至开仓放粮赈济灾民……

这些事情倒是真的，老刀子倒也听说过不少。

在十分话语中只有一分提及问道途，是吴刚的策略——而这一策略自然也是卓有成效的。

潜移默化间，老刀子看问题的角度和深度，悄然发生了变化。

他不再仅仅盯着眼前的矿石和工钱，开始下意识地思考丹鼎阁定价的不公根源，思考紫阳宗征调政策背后的冷酷算计。

更重要的是，他作为南区矿工中有威望头目的身份，使他的变化和偶尔透露出的只言片语，开始产生影响。

他依然不可能公开说什么——实际上对问道途的印象虽然好了些，但也确实没太好。

但在矿洞休息的间隙，在劣质酒馆的一角，当手下兄弟或相熟的散修抱怨“这日子没法过了”、“上面心太黑”、“后生们没出路”时，老刀子闷头喝一口酒，可能会沉沉地叹口气，无意间接上一句：

“听说云霞山那边，散修死了，抚恤是足额发的，家里人也有人安置。”

或者，在有人羡慕某个依附家族子弟又得了好处时，他擦擦脸上的灰，冷不丁冒出一句：“依靠别人赏饭，终究不稳当。我好像听说，云霞山那个叶先生，传道不怎么看重出身和灵石，有个叫陈默的，以前也就是个跑单帮的，现在可不一样了。”

他甚至连表演的痕迹都没有——因为根本没有表演，全都是真正拉闲散闷的口气。

配合着他本身硬气、实在的形象，反而比任何激昂的演说更具说服力。

听者会愣一下，然后追问：“真的假的？老刀子你听谁说的？”

老刀子通常只会摇摇头：“道听途说，谁知道呢。”

但那种子，已经借着他自己都没意识到有多不经意的话语，种进了听者的心里。

他们会私下讨论，会去打听，会结合自身遭遇去想象——“如果那是真的……”。

\vspace{2em}

吴刚塑造的，不止一个“老刀子”。

“老刀子”们或许不清楚背后运作的万仙盟网络全貌，甚至不清楚吴刚具体是个什么东西，也不知道哪个“到处乱窜的女侠”是何许人也。

但他们能模糊地感知到，有一群不同于守旧联盟的、带着些许温度的家伙——甚至有时候自己都不理解这些。

于是，这些人便成了吴刚网络中最基层、也最稳固的无形节点。

他们甚至不会打出旗号，不会聚众闹事；但在日常生活的每一个缝隙里，他们偶尔泄露出的一两句“牢骚”或“感慨”，持续地瓦解着对守旧联盟的敬畏与顺从，播撒着对另一种可能的好奇。

而陈默，这个活生生的“榜样”，则像巡回展示的珍宝，在绝对安全的掩护下，被吴刚带到几个类似黑石集的关键节点，与这些“节点”人物进行更深入的接触。

在伪装过的密室中，陈默没有任何夸夸其谈。他只是平静地讲述自己作为一个普通散修曾经的窘迫与迷茫，讲述在云隐别院第一次听到叶牧谦讲解“仰不愧于天，俯不怍于人”时的震撼，讲述转修问道途后，那种虽然缓慢却扎实的进步，以及内心逐渐明晰的笃定感。

他展示的是一种迥异于寻常散修的精气神——沉稳、内敛，眼神清澈而坚定，没有长期被盘剥压榨者的畏缩与怨毒，也没有骤然得势者的骄狂。

老刀子这样的人，看人极准。他们能看出陈默不是装的，那种从骨子里透出的踏实与尊严，做不得假。

陈默的存在本身，就是最有力的证据，证明了“问道途”和云隐别院，至少为他们这样的人，提供了另一种活法——或许依然艰难，但更有尊严，更有希望。

“陈默道友，”老刀子在一次秘密会面后，忍不住低声问，“叶先生……真如传闻所说，传道不论出身，不索贿赂？”

陈默认真点头：“老师有教无类。在别院，一切用度、功法，皆按需分配，按劳、按功获取。欺压同门、盘剥弱者，是首要大忌。”

老刀子沉默良久，重重吐出一口浊气。

“妈的，这问道途还真有点东西。”

\vspace{2em}

于是，发动群众的效果，开始以另一种更绵长、更深刻的方式，在黑石集以及类似的土壤上显现出来。

效果首先体现在最直接的经济层面——这自然是最快的，拜此地散修的实际领袖、与吴刚秘密签订条约的老黑石所赐。

丹鼎阁又派人来收购药材矿石了。这个精明的管事来这里也不是一回两回。却发现这个月黑石集上交的配额，虽然面子上勉强够数，但品质上佳的“尖货”比例明显下降，且总数量竟然比预估的“正常产出”少了那么一成半成。

“怎么回事，老黑石，你们这些家伙是不是偷懒了，没好好替我们收货？”

老黑石一脸愁苦地摊手：“仙师明鉴啊！今年矿脉贫瘠，妖兽闹得也凶，兄弟们实在是拼了老命了……”

“抗命是要杀头的。”管事无奈地道，“我也是有命令在身的，收货物收不上来我也不好做啊。”

老黑石也愁眉苦脸地递上来一张纸。

“您看，伤亡名单都在这呢。”

管事也没什么可说的，毕竟都到这种程度了，死者为大；况且收货的面子上也够数了，他对上级也有个交代。

而与此同时，吴刚的秘密渠道，却总能以比官价高出三到五成的“公道价”，稳定收到一批成色不错的货物。散修们拿到实实在在的更多报酬，对老黑石和那神秘的“外路朋友”自然心存感激，心照不宣。

\vspace{2em}

其次体现在对联盟统治的消极反抗上——这个自然要慢得多了，毕竟改变人心是一个更加缓慢的过程。

当紫阳宗的传令修士再次来到黑石集，要求紧急征调三十名灵枢境散修去填补前线某个危险区域的防务缺口时，响应者寥寥。

以往靠利诱和威逼总能凑齐的人数，这次磨蹭了三天，也只来了十几个老弱或实在走投无路之人。

传令员自然是不信老黑石等人“动员不上来”的说辞的，他倒也务实，选择亲自下去动员散修。

老黑石同样跑前跑后，跟着传令员“苦口婆心”动员，却挨了老刀子一通臭批：

“驴*的紫阳宗的狗东西，你先还我侄儿再说！老黑石你他妈的也是，怎么还敢来这，再敢来我把你揍出去！”

一连几处，皆是如此。老黑石和传令员也是挨了不少唾沫星子。

最终，老黑石也只能无奈地回禀：“仙师，您也看到了，人心涣散，惧怕前线死地，小老儿实在是尽力了……”

传令员灰头土脸，挨了一天骂饶是谁也都不好受，也只能带着这十几个人回去复命了。

这自然也仅仅是个缩影，但随着时间流逝，征调令的执行效率确实在大打折扣。

联盟在黑石集及周边区域的驻军换防时间、巡逻路线、仓库守备的薄弱时段……这些信息，开始通过那张地下网络，悄然流向吴刚，经过汇总分析后，又变为更有效的行动指导，流向沈瑶等人。

而守旧联盟对黑石集等地内部的真实人心动向、暗流涌动，却如聋似瞎，只知道收税、招人不太好办了。

\vspace{2em}

最后，也是最重要的，是一种人心氛围的微妙转变。

集市上，茶摊酒肆里，公开咒骂联盟的声音或许还不敢有；但那种麻木的顺从在减少，取而代之的是一种闪烁的、心知肚明的沉默，以及私下交流时眼神中蕴含的别样意味。

当有联盟修士趾高气扬地走过时，背后投去的目光里，敬畏在减少，隐忍的厌恶与一种近乎怜悯的疏离感在增加。

那句“仰不愧于天，俯不怍于人”，开始像一道微弱却执拗的光，照进一些不甘沉沦的心灵深处。

吴刚领导的这张网，如同生长在守旧联盟庞大躯体深处的无形菌丝。它不发动正面的攻击，不攫取显赫的城池，却以惊人的耐心与韧性，无声地缠绕、渗透进联盟基层统治最细微的根系之间。

它缓慢而持续地吸走本该属于联盟的养分，悄然松动其统治的土壤，并将另一种关于生存方式、关于尊严与未来的可能性，如同孢子般，悄无声息地播种在无数散修的心田。

\vspace{2em}

当然，暗流涌动之下，必有波澜。

守旧联盟能做到现在这种大小，如果其说其中枢是蠢物那有些太不负责任了；只是，其臃肿的躯体对发生在末梢神经的细微刺痛，反应往往迟钝而傲慢。

这件事一开始甚至没有捅到陈炎那里，

紫阳宗前线那位负责后勤协调的中级执事，翻看着下面陆续呈上来的卷宗，眉头拧成了疙瘩——这个月从黑石集等几个固定供应点收上来的低阶矿石和药材，平均品质比上月低了半成，总量也差了那么一丝。

呈报上来的理由千篇一律：矿脉贫瘠，妖兽滋扰，人手折损。

“废物！”执事低声骂了一句，随手将卷宗丢到一边。

在他眼中，这只是下面那些贪婪又无能的附庸小家族或散修头目惯用的伎俩，想借此讨要些补贴或降低配额。

相较于千头万绪的战报和日常管理，这等小事尚不足以惊动更高层。

要是这种鸡毛蒜皮都上报，那免不了挨一顿批评：“此等小事还拿来犯我，你能力是不是有点过于有限了？”

是啊，这些下级，能力是不是过于有限了？

这个执事有些恼火，批了一句：“严加督促，务必足额，否则以延误军机论处！”便将这份报告归档了。

几乎同时，另一份来自某个地方驻军的报告也送到了他的案头，报告称辖区内“流言频起”，有低贱散修私下议论前线战事不利，甚至妄言紫阳宗苛待仆从军。

执事的火气更大了，提笔批复：“妖言惑众者，一经查实，立斩！驻军加强巡防，以儆效尤。”

他依然没有将这几件“小事”联系起来。

在他，乃至守旧联盟大多数中上层修士的思维定式里，后方的些许杂音，无非是管理懈怠或愚民无知，派几队更凶悍的巡逻队下去敲打一番，杀几只鸡，猴子们自然就老实了。

他们尚未意识到，这早就不是简单的懈怠或愚昧了；而是一种无声的、织网般蔓延的抵抗。

真正让“小事”开始引起区域层面注意的，是沈瑶那边愈演愈烈的打砸抢。

距离黑石集数百里外，一处由丹鼎阁某个附属家族控制的灵谷转运站，在某个深夜燃起了冲天大火。不仅转运站本身被焚毁、数千担灵谷不翼而飞，更有数名家族修士被当场格杀。

现场留下用血迹涂抹的大字：“取不义之粮，还于饥民！”

只知道那名——或者那伙——“贼人”擅长短刀，轨迹极为精确，基本上全奔着一刀毙命去的。

几乎每隔十天半月，类似的事件就会在守旧联盟控制区域的薄弱地带上演。

有时是押送税款的队伍在半道被劫，护卫被杀，灵石不翼而飞；有时是某个盘剥甚重的小家族仓库被撬开，里面的物资一夜之间消失大半，只留下些散给附近穷苦人家的痕迹。

行事者来去如风，手段狠辣精准，绝不纠缠，一击即走。渐渐地，一个模糊的代号在底层流传开——“那伙煞星”，或者更具体点，“那个穿蓝衣服、刀快得像鬼一样的女人”。

\vspace{2em}

于是这些事件终于不再是“小事”。

损失实实在在，影响恶劣，更重要的是，它们指向了一个明确的、有组织的势力活动。

不知道这贼人是哪来的，但是有贼人到处活动这事情是真的。

几位主事者们坐不住了。

一次紧急议事中，负责本区域防务的紫阳宗执事面色不善地敲着桌子：“不能再姑息了！先是物资短少，征调不力；现在匪患公然劫掠，袭杀我修士！这已不是简单的散修和老百姓了，必须要出重拳！”

另一位丹鼎阁的代表则更关心实际利益，阴恻恻地说：“尤其是黑石集一带，近来颇有古怪。收缴上来的货物品相持续走低，必有猫腻。说不定，那里就是匪徒销赃、补充给养的窝点之一！”

他们迅速达成了共识，双管齐下。

一方面，调集精锐，加大对那支“流窜匪帮”的清剿力度，重金悬赏其人头；另一方面，对黑石集这类“可疑区域”施加高压，派加强的巡逻队进驻，明查暗访，务必揪出内鬼，掐断任何可能与“匪帮”联系的渠道，重新树立联盟的绝对权威。

\vspace{2em}

数日后，第一支背负着震慑与清查双重命令的联盟巡逻队，开进了黑石集。

这支队伍由二十名修士组成，远超平日巡防的规模，其中领队是一名玄脉境通络期的紫阳宗资深弟子，面色冷峻，眼神里带着不加掩饰的审视与倨傲。队伍中还有两名来自丹鼎阁、专管财务的执事，专门负责查验货物账目。

集市主干道上，人群下意识地分开，原本嘈杂的讨价还价声低了下去，无数道目光或畏惧、或麻木、或带着隐晦的疏离，从街道两旁、摊位后面投来。

巡逻队队长很满意这种效果，他要的就是这种令人窒息的威压感。

他们首先直奔老黑石的“公事房”——一间简陋的石屋。老黑石早已得到风声，恭敬地迎了出来，腰弯得很低，脸上堆满了讨好的、属于小人物的愁苦笑容。

“仙师们大驾光临，有失远迎，恕罪恕罪！”老黑石搓着手，语气卑微，“可是上峰有什么新指令？小老儿一定尽心竭力去办！”

队长没理会他的客套，径直走入石屋，目光如鹰隼般扫过每一个角落。两名丹鼎阁执事则毫不客气地开始索取近三个月的出货账册、矿工伤亡记录、以及所有与丹鼎阁交易的明细凭证。

老黑石早有准备。账册记得清清楚楚，明面上的出货量、品质分类完全符合“官方要求”，只是旁边用朱笔小字备注的“矿脉贫瘠预估减产”、“妖兽袭击损失”、“新矿工伤亡抚恤支出”等条目，密密麻麻，触目惊心。

伤亡名单更是厚厚一摞，每一个名字后面都按着手印或简单的标记，透着一股沉甸甸的、令人不悦的真实感。

一名执事皱着眉头翻看，试图找出涂改或作假的痕迹，却一无所获。他抬起头，狐疑地盯着老黑石：“老黑石，这数目，可经得起查验？若是虚报伤亡，侵吞物资，你知道后果。”

老黑石脸上的愁苦简直要滴出水来，他指着账册，声音都带上了哽咽：“仙师明鉴啊！借小老儿一百个胆子也不敢啊！实在是……唉，仙师们若是不信，大可去矿洞看看，去那些死了男人的家里问问！这日子，难啊！”

他表演得如此自然，将一个小人物在夹缝中求生存、既要完成任务又无力回天的形象刻画得入木三分。至于这些账目，在吴刚等人的参与和指导下，七分真、三分假，假的地方大致上也算天衣无缝。

巡逻队队长哼了一声，对这种诉苦把戏不感兴趣。

他手一挥，两名丹鼎阁财务执事立即开始上前查账，算盘打得噼里啪啦。

当然，巡逻队队长更关心的是“匪患”线索。

“最近集市上，可有什么生面孔？有没有人私下打听云霞山，或者交易来路不明的物资？”队长冷声问道，目光锐利如刀，仿佛要剖开老黑石的脑子看看。

老黑石茫然地眨眨眼，做努力思索状：“生面孔……这黑石集哪天不来几十号生面孔？都是讨生活的苦哈哈。云霞山？哎哟，那可不敢提，那是叛逆之地啊！至于来路不明的货……”

他压低了声音，显得小心翼翼，“仙师，咱们这儿规矩，来历不明的东西不收，怕惹麻烦。倒是……倒是偶尔听说，有些胆大包天的家伙，会偷偷把一点私货卖给那些行踪不定的游商，价格能高些，但风险也大，小老儿是严令禁止的！”

他把“偷偷”和“游商”咬得很重，既暗示了可能存在的地下交易——这符合联盟对平民百姓、底层散修的刻板印象，又把自己撇得干干净净。

与此同时，巡逻队的其他成员已经分散开，在集市上进行抽查。他们闯入几家较大的药材铺和矿石行，翻箱倒柜，盘问掌柜和伙计，语气凶狠。

集市上一时间鸡飞狗跳，人心惶惶。

然而，他们很快发现，面对的是一个无比配合却又无比空洞的群落。

所有的家伙此刻都和常规没有两样。老刀子秦厉正光着膀子，和一群矿工蹲在街角啃着干粮，满脸灰尘，眼神和其他人一样警惕而麻木。

巡逻队员过来盘问，他夹枪带棒地骂了两句，表示什么都不知道，只念叨着工钱难挣，饭不够吃，“最好少收点税”。

所有的关键交易，早已转入更深、更隐秘的渠道。

吴刚的秘密运输线采用了更灵活的单线联系和分段接力，物资在联盟巡逻队到来前，就已通过不同的小路，化整为零地转移或藏匿起来；然后再分批混在更多的货物中，或采取其余手段，把这些“灰色”货物逐渐地变成经得起查验的好货。

通俗来说，那叫“洗钱”。

洗钱归洗钱，至少他这个洗钱确实起到了积极效果。

现在集市上流通的，都是经得起查验的清白货物。

巡逻队折腾了大半天，近乎一无所获；查账的两个财务人员也没看出什么蹊跷。

但他们想找的内鬼、联络点、大批违禁物资，一样也没发现。

领队的队长脸色越来越难看。他凭直觉感到不对劲，这里的气氛太过正常，正常得过了头。

可他抓不到把柄。

最终，他只能抓了几个平日里就小偷小摸、或是酒后胡言乱语说了几句怪话的倒霉蛋交差。

临行前，巡逻队队长严厉警告老黑石和集市上几个有头脸的人，必须严密监视可疑人等，及时上报，否则以通敌论处。

然后，带着一丝悻悻然和犹疑，巡逻队离开了黑石集，前往下一个“可疑”地点。

\vspace{2em}

一连几处，皆是如此。

\vspace{2em}

然而，当后方“物资损耗异常”、“征调不力”、“流言频起”乃至“匪患劫掠”的报告太多、无法处理，最终如同雪片般从各个方向汇集，甚至有几份直接呈送到了坐镇前线的陈炎案头时，事情的性质就变了。

尽管陈炎的主要精力仍被云霞山这个“硬骨头”牵制，但后方此起彼伏的烽烟，尤其是那些精准打击补给线、动摇统治基础的行动，让他无法再将其视为简单的“流窜匪患”或底层散修的抗命。

一道带着陈炎冰冷意志的命令被迅速下达：后方各区域，必须严加整肃，加大巡查与清剿力度，务必在最短时间内扑灭这些“星星之火”，保障后方安稳与前线的物资供给。

压力如同不断收紧的绞索，开始套向黑石集这样的地方，也迫使沈瑶和吴刚的行动更加艰难和险象环生。

双方的探子、小股部队在城镇外的荒野、山林间的遭遇变得频繁。

在一次对某处小型联盟物资中转站的侦察任务后，沈瑶带着游击队的三名精锐队员，趁着夜色撤离，准备返回位于深山中的临时营地。

他们行动如狸猫，借着山林阴影和沈瑶以千幻流光诀制造的细微视觉扭曲，几乎与夜幕融为一体。

就在穿过一片稀疏的林地，即将进入更茂密的山岭时，前方突然传来杂乱的脚步声和低沉的呼喝声，还有几道神识之光如同探照灯般扫过林间空地——是一支约十五人的联盟巡逻队，似乎是在执行例行的夜间巡防，领头的是一名神色警惕的玄脉境修士。

双方距离不过百丈，且处于相对开阔地带，遁走极易引发灵力波动而被察觉。

硬闯？对方人数占优，一旦纠缠，附近很可能还有其他巡逻队。

现在还没被发现，还有时间伪装！

电光石火间，沈瑶眼神一凛，手势微动。

队员们心领神会。

只见四人几乎同时，发生了细微却显著的变化：沈瑶身上那属于精锐修士的凌厉气息瞬间收敛殆尽，换上的是一种长途跋涉后、饱经风霜的疲惫与谨慎，她甚至下意识地弯了弯腰，仿佛背负重物；其他三名队员则迅速将身上的衣装弄得凌乱，脸上也抹上些尘土，眼神变得茫然又带着点小民的畏缩。

一切准备完成，几人调整方向，略显慌乱地朝着巡逻队即将经过的方向“撞”了过去。

常言道：最危险的地方，就是最安全的地方。

“什么人？！站住！”巡逻队立刻发现他们，数道神识和兵刃瞬间锁定。

沈瑶立刻举起双手，用带着浓重口音、有些颤抖的方言喊道：

“军、军爷饶命！俺们是北边野猪岭的采药人，进山寻点山货，迷了路……这、这才转出来……”

领头的玄脉境修士眉头紧皱，神识仔细扫过四人。

修为最高者不过灵枢中期，气息虚浮杂乱，衣着普通甚至寒酸，身上带着淡淡的草药泥土味和汗味，确实像极了常年混迹山野的底层采药散修。

他目光锐利地扫过沈瑶的脸——那是一张粗糙、黝黑、带着惶恐的农妇面孔，毫无特点。

但野猪岭离这边有五十里地！

“野猪岭？跑到这五十里外的林子来采药？当老子好糊弄？！”

巡逻队长厉声喝道，同时示意手下散开，形成半包围。

沈瑶身体抖得更厉害，带着哭腔道：“军爷明鉴啊！俺们……俺们也是没办法！野猪岭那边被征了‘入山税’，俺们交不起，听说这边林子深，没人管……就、就冒险过来了……谁知道这鬼林子这么大……”

她一边说，一边从怀里颤巍巍掏出一个破旧的布袋，倒出几株品相普通、还沾着泥土的“灰线草”和“地根藤”，都是黑风林附近常见的低阶药草。

“军爷，这是俺们今天采的……都、都孝敬军爷，只求军爷指条明路，放俺们回去吧……”

巡逻队长看着那几株再普通不过的药草，又看了看眼前这四个采药人瑟缩恐惧的模样，心中的疑窦去了大半。

这种为了逃避苛捐杂税、冒险进入陌生区域讨生活的散修他见多了，都是些穷得叮当响、修为低微的可怜虫。

抓起来压榨，也榨不出什么油水。

更何况那农妇容貌根本入不了他的脸。

他嫌恶地挥挥手，像驱赶苍蝇一样：“滚！以后不许再靠近这片区域！再让老子看见，按奸细论处！”

沈瑶四人如蒙大赦，连连鞠躬道谢，然后互相搀扶着，脚步踉跄却速度不慢地“逃”向了与真实营地相反的方向，迅速消失在黑暗的林地中。

直到彻底脱离对方神识范围，四人才停下，互望一眼，都从对方眼中看到一丝紧绷后的放松。

沈瑶面容一阵水波般的模糊，恢复了原本的模样。

这次蒙混过关，靠的是对底层生态的精准模仿、毫无破绽的伪装，以及沈瑶那足以完美控制自身气息甚至在一定程度上干扰低阶修士感知的千幻流光诀；但第二次，他们就没有这么幸运了。

\vspace{2em}

真正的危机，来自内部的背叛。

吴刚发展的地下网络，不可能永远铁板一块；在巨大的压力、或许还有守旧联盟暗中开出的价码诱惑下，人性中的脆弱与贪婪总会暴露。

一个代号“灰鼠”、负责传递情报和转运少量物资的下线，被当地一个依附丹鼎阁的赵家秘密策反了。

赵家老祖是真符境摹形期的修士，故而在本地势力颇大；早就对境内“不安定因素”深感恼火，得到了上峰“不惜代价肃清”的严令。

毕竟人都是有软肋的；有软肋，那就可以拿捏。

他们抓住了灰鼠一个至亲的把柄，并以重利许诺，逼其就范。

威逼利诱之下，灰鼠最终还是选择叛变。

或许是因为他想把队友分期卖个更高的价钱，也有可能是因为他良知尚存，灰鼠没有立刻暴露全部网络，而是选择向赵家提供了一条他认为是“大鱼”的情报：

三天后，游击队将秘密潜入镇外废弃的老矿洞区域，接收一批由吴刚网络从其他地方转运来的、较为珍贵的疗伤丹药和一批用于制造简易阵法的材料，并计划在那里休整一夜！

他赌的是用沈瑶这条“大鱼”，来换取自己和家人的安全，以及赵家许诺的荣华富贵。

情报被火速呈报。赵家老祖亲自出马，协调了当地紫阳宗驻军，抽调精锐，秘密在老矿洞周边布下了天罗地网。

他们甚至动用了能够屏蔽低阶传讯、干扰灵力感应的阵旗，只待沈瑶小队进入，便一网打尽。

沈瑶对此一无所知。

灰鼠这条线之前传递过几次无误的情报，且这次接头地点和时间，是通过吴刚网络内部相对安全的渠道确认的，应当无误。

直到他们深入矿洞区域腹地，准备发出接头的暗号时，异变陡生！

“轰！”

四周山壁上，预先埋设的爆裂符骤然激发，制造了相当的巨响和混乱，打断任何可能的遁术起手。

与此同时，至少超过五十道身影从废弃的矿坑、堆积的废石堆后、甚至是从地下破土而出！

居中那位白发老者，周身符光隐隐，身前悬浮着一枚不断变幻形状的土黄色符印，正是赵家那位真符境摹形期的老祖！

其左右，则是两名玄脉境后期的赵家长老；外围，是数十名紫阳宗驻军和赵家护卫组成的包围圈，刀剑出鞘，弓弩上弦，灵光锁定了中心的沈瑶小队。

“你就是沈瑶？好个云霞山余孽！束手就擒，否则今日此地，便是你的葬身之处！”

赵家老祖声如洪钟，带着志在必得的杀意。

中计了！

沈瑶瞬间明白过来，眼神冰冷如万载寒冰。她自认凭千幻流光领域以及堪比真符境的战斗力，这些人估计是留不下她本人的。但是问题是，如果赵家老祖故意放过她，选择强行留下其他游击队员……

她没有任何废话，清叱一声：“跟我冲！”

话音未落，她已身化一道变幻不定的流光，直扑包围圈西北角。

根据她瞬间的观察，那边地形复杂，有突围的可能。

千幻流光诀全力催动，她身影仿佛同时出现在数个方位，道道虚实难辨的刀光泼洒而出，瞬间将那片区域的守军搅得人仰马翻。

队员们也红了眼，知道已是绝境，唯有血战求生。

他们结成小型战阵，不顾自身防御，将攻击催发到极致，紧紧跟随着沈瑶撕开的缺口。

“想走？！”赵家老祖冷哼一声，身前土黄色符印光华大盛，瞬间化作一面巨大的、布满玄奥纹路的岩石屏障，轰然砸落在沈瑶的突围路径上，更有一股沉重的重力场弥漫开来，让沈瑶和队员们的动作骤然一沉。同时，他袖中飞出数道金光，竟是成套的飞剑符宝，凌厉无比地攒射向沈瑶。

另外两名玄脉境长老也联手攻来，术法光芒呼啸。

“靠，你们先走，真符我能拖住！”

沈瑶骂了一声，立即旋身下达指令。

但没有人听从命令。

“队长，我们跟您一起！”

“你们是笨蛋吗！”沈瑶急了。

但脑子长在别人身上。

很快，其中一名游击队员被飞剑符宝穿透胸膛，当场壮烈牺牲，眼见得救不活了。

另一名队员在引爆身上携带的所有爆炎符，暂时逼退一片敌人后，也被乱刀砍倒。

“要活的，你们干什么呢！”

赵家老祖狂吼道。

很快，鲜血染红了废弃的矿场。沈瑶自己手臂也被一道符宝金光擦过，受了些伤。

她咬紧牙关，知道再拖下去，全体都要葬送在此。

“散开！各自突围！老地方汇合！”她厉声喝道，同时拼着受了赵家老祖一记符印重击，虽被震伤，却借力向后急退。

忽然——

“弥焰斩！”

众多追兵忽然听得一声男子暴喝，却听不真切是哪里发出来的。

“什么，那个姓白的也来了？”

“那个弥焰斩可是擦着边就死、碰了些就亡啊！”

赵家老祖也忌惮起来，毕竟白言冰加上沈瑶，两个真符……

就在这众人因犹疑而迟滞的一瞬！

“轰隆——！！”

震耳欲聋的爆炸声中，电蛇乱窜，火光冲天，剧烈的黑紫剑气（？）风暴暂时遮蔽了视线和感知。

随即，是浓到化不开的烟尘。

沈瑶身影在爆炸的掩护下，左右手上各抓了一个受了伤、无法全力飞遁的家伙。随即，身形如鬼魅般连连闪烁，千幻流光诀催动到极致，气息变幻不定，甚至模仿出数道向不同方向逃窜的虚影。

赵家老祖第一个反应过来。

这黑紫剑气徒有虚表，毫无杀伤力啊！

还能这么打架的？

怒吼着驱散烟尘，神识疯狂扫荡，却只抓到几缕迅速消散的幻影和远处隐约的灵力波动，根本无法锁定沈瑶真身。

他脸色铁青，知道最关键的“大鱼”脱钩了。留下的两个游击队员也已经成了尸体，盘问算是盘问不出来了。

“娘西匹，给我追！”

——上哪去追啊！

\vspace{2em}

沈瑶全力飞遁数十里，才终于停下，算是勉强逃出了包围圈。

跟随她成功突围的，只有三名玄脉境巅峰的队员，都受了些轻伤；以及另外两名被沈瑶抓出来、身受重伤的队员，但算是彻底失去了战斗力。

“对了，白师兄哪里去了？”一名帮助沈瑶临时为重伤号包扎伤口的队员随口问道。

“白师兄？哪个白师兄？”沈瑶疑惑道。

“白言冰师兄啊！”

沈瑶白了他一眼，“那是我用幻术领域模仿的！你当他真能来啊，那山上防务谁负责？”

队员面色有些变了：没想到沈瑶的幻术竟然到了这种地步，甚至连白言冰的杀招——弥焰斩都能轻易模仿！

沈瑶看出他的惊奇，又补充道：“那黑紫剑气也是假的！”

队员哑然，笑也不是，不笑也不是。

沈瑶逃出生天、立即通过紧急渠道将“任务失败”的消息传递给吴刚。

吴刚暴怒，动用地下网络的力量查明事实后，在灰鼠试图领取赵家赏金的前夜，将其秘密处决。

经此一役，“沈瑶”这个名字及其大致特征，以及其余几个逃出生天的队员，被守旧联盟正式列为高度危险的通缉要犯，画像与悬赏令贴满了后方许多城镇的关口要道。

敌人已经深刻认识到她的威胁，并会调动更多资源、采取更狡猾的手段来对付她和她的队伍。

当然，这对沈瑶而言几乎毫无意义——千幻流光诀一开，她想换成什么模样，就是什么模样。

通缉令上那张脸，在内域大乱结束之前，估计永远也不会再次出现了。

单纯的通缉自然是无意义的，联盟随即在内域全域大索，大力镇压可能的反对派活动，搞得整个内域鸡犬不宁，民不聊生。

然而，高压往往带来更强烈的反弹。越是镇压，那些分到过粮食的贫民、那些受过盘剥的散修、那些偷偷修炼《养气篇》感受到体内微弱元气增长的年轻人，心中对守旧联盟的憎恶和对问道途的向往就越是强烈。

地下网络、游击队伍在一次次的打击中变得更加隐秘和坚韧，如同野草，火烧不尽。

\vspace{2em}

很快，消息通过秘密渠道断断续续传回云霞山。

尽管山上的防御战依然惨烈，物资依然紧缺，但当叶牧谦将沈瑶、吴刚他们在敌后发动群众、分发物资、传播理念的事迹选择性地告知守军时，一种难以言喻的振奋在山间弥漫开来。

他们不是在孤军奋战！

他们的抗争，正在遥远的地方激起回响，正在动摇敌人看似稳固的统治基础！

这种认知，极大地鼓舞了防守者的士气。

白言冰在指挥防御时，有时会吼出：“想想沈瑶师妹他们在敌后为民请命！我们守住这里，就是守住希望的火种！”

陈千羽在救治伤员时，也会轻声鼓励：“吴刚师兄他们找到了更多药材来源，大家坚持住，我们的药田又快丰收了。”

而叶牧谦收到密信时，更是久久不能平静。

他的思路，是正确的。

实践也检验了真理。

本想提笔回复一些鼓励的话，但最终千言万语，终究不如井冈山上那光耀千古的八个字：

“星星之火，可以燎原。”

山上山下，尽管被重兵隔断，却通过精神的共鸣和战略的联动，形成了一个坚韧的整体。

云霞山如同吸引雷霆的孤峰，承受着正面最大的压力；而敌后战场则如地下奔涌的暗流，不断侵蚀着敌人的根基。

当陈炎再次于中军大帐中接到后方接二连三关于“匪患”“邪说流传”的急报时，他的脸色终于不再是单纯的阴沉，而是带上了一丝不易察觉的焦躁。

他发现自己面对的，不再只是一个龟缩防守的军事据点，而是一个正在用他完全不能理解的方式，将战争扩展到整个社会层面的可怕对手。

“叶牧谦，你到底是什么人……？”

陈炎不禁冷汗淋漓起来……

\vspace{2em}

沈瑶在许多个“灰谷镇”斩开的仓门后纷飞的粮食，吴刚在暗室中与许多个“老黑石”们推心置腹勾勒出的利益与尊严的蓝图，成了守旧联盟大后方的星星之火。

看似微弱，却因其内核燃烧的是“公道”与“希望”的燃料，而拥有了顽强的生命力。它们开始在守旧联盟统治的厚重阴影下，倔强地闪烁、蔓延。

最直观的对比，发生在每一天，每一处。

守旧联盟的大军或其爪牙过境，无论面对的是名义上的附属城镇还是自家的后方腹地，“征发”是唯一的逻辑。

粮食、药材、矿石，乃至低阶修士的性命，都在“道争大业”的冠冕堂皇之名下，被理所当然地索取。

尚有良知、愿意通融的低阶军官固然存在，但执行这道命令的更多低阶军官、税吏、附属家族的狗腿子们，将这视为大发横财、作威作福的良机。

他们敲骨吸髓，中饱私囊，视散修与凡人为无血无泪、可以无限榨取的资源矿藏，其行径之暴虐贪婪，较之山野匪类有过之而无不及。

失去儿子的母亲眼泪未干，夺走口粮的差役鞭影又至。

与此形成刺眼到令人心颤对比的，是沈瑶、吴刚小队那秋毫无犯，开仓济民的作风。

许许多多个“灰谷镇”的故事，早已被口耳相传，添加上百姓最朴素的想象与期盼，演化成数个略有不同的版本，在周边的乡镇、矿场、散修聚集地悄悄流淌。

起初，人们只当是乱世中一个慰藉心灵的奇谈，带着深深的怀疑与本能的观望；然而，现实是最好的教员。

当守旧联盟的征粮队因“损失”而将亏空变本加厉摊派到邻近村镇时，当催缴“特别军税”的皮鞭再次抽打在早已家徒四壁的农户背上时，那份被强力压制的怨愤，如同被堵塞的岩浆，开始炽热地寻找每一个可能的裂隙。

第一个大胆的散修，在月黑风高的夜晚，怀揣着难以抑制的愤怒与一丝微茫的希望，循着模糊的传言，冒险找到了沈瑶小队曾经短暂歇脚、现已废弃的山洞。

他留下的不是金银，而是一张用炭笔画在粗麻布上的简陋草图，上面歪歪扭扭却清晰地标注了一支征粮队返回驻地的必经之路、人数、以及押运修士的大致修为。

这份冒着生命危险送出的情报，成了点燃更多火焰的火种。

沈瑶——这回谨慎得多了——根据情报，多次踩点之后精准设伏，不仅截获了粮队，更再次将粮食分发给沿途受尽盘剥的村庄。

观望，在这一刻，化为了实实在在、浸透着泪水的感激。

人心的天平，开始了无声却致命的倾斜。

支持的形式，从此如雨后山林里的菌菇，在仇恨与渴望的潮湿土壤中，悄然萌发，形态万千，生命力顽强。

有的支持，如同那位无名散修，是关键时刻递出的一张纸：一张写着某条补给线护卫换班间隙的纸条，塞进赶集农妇的菜篮底；几句关于某个矿点监工今日醉酒、守卫松懈的闲谈，“无意”飘进酒馆角落里看似打盹的旅人耳中。

这些零碎的情报，通过吴刚再次精心编织、已愈发细密坚韧的地下网络，被迅速筛选、拼接、传递。

最终，它们可能化为沈瑶“千幻流光”下的一次精准斩首，救出几名即将被折磨至死的矿奴；或是引发某个偏远税卡守卫的“哗变”，整队人马带着武器和税款消失在山林，据说是投奔了游击队。

\vspace{2em}

寒鸦岭，这名字听着就透着一股子穷苦与荒僻。

岭下散落着几十户人家，多是土坯垒墙、茅草覆顶，日子过得比岭上的石头还硬。张老栓家就在村西头最边上，两间矮趴趴的茅屋，像是随时会被山风吹散架。

约莫三天前的后半夜，张老栓被屋后柴垛一阵窸窸窣窣的动静惊醒，还夹杂着压抑的、粗重的喘息。他抄起门边的柴刀，壮着胆子摸黑出去，月光下，只见柴垛旁歪着个人，半边身子都被暗红色的血浸透了，脸色白得像死人，只有胸口还在微弱地起伏。

那人穿着看不出本色的粗布衣服，但手里还死死攥着一把卷了刃的刀，刀柄缠着的布条都被血粘住了。

张老栓的心提到了嗓子眼。

他认得这打扮，近来山外边不太平，传得沸沸扬扬，说有一伙专跟“老爷们”作对的好汉，领头的是个女侠，来去如风。

眼前这位，八成就是那伙人里的，看样子是在别处吃了亏，一路逃命躲到了这山旮旯。

救，还是不救？

张老栓蹲在黑影里，嘴里的旱烟吧嗒吧嗒，却半天没吸进一口。救了吧，万一被王老爷或者上面下来的兵爷知道，这就是包庇匪类，要杀头的。

不救吧……这后生眼看就剩一口气了，扔在外头，不是冻死就是被野狗啃了。

他想起了上个月，岭东头老赵家。赵家的独苗后生被王老爷的征粮队硬拉去当了脚夫，说是去前线运粮，结果人没回来，只捎回句话，说是“遭了流匪，不幸没了”。

而王老爷的管家来“抚慰”时，只丢下两斗谷子，那架势，跟打发叫花子没两样。

一个壮劳力的性命就顶两斗谷子！

老赵头当时就呕出口血，没挺过三天。

他又想起了大概半个月前，村里悄悄流传的消息。说百里外的“灰谷镇”，老爷家那个屯了满仓粮食、却年年加租的大粮仓，被一伙人给砸开了，里头的粮食全分给了镇上的穷户。

消息传得有鼻子有眼，还说那伙人只抢老爷的，不碰穷人家一粒米。

当时村里人听了，只当是故事，心里却都觉得解气，晚上睡觉都踏实了些。

柴垛边那伤员的呼吸更微弱了。

张老栓把烟锅子在鞋底磕了磕，站起身，像是下定了决心。

他左右瞅瞅，确认四下无人，这才费劲地把那人连拖带拽，弄进了屋。家里空荡，藏不住人，他一眼瞅见了屋角那堆平时做饭用的柴火。老伴去得早，儿子也被征走没了音讯，就他一个孤老头子，柴火用得少，堆得老高。

他挪开面上干燥的柴捆，露出后面墙壁。这墙是土坯的，早年防山匪，他在后面偷偷挖了个小窑洞，不大，也就刚好能蜷个人，洞口拿块破木板遮着，再堆上柴，外人绝看不出来。洞里阴冷潮湿，常年不见光，有一股子霉土味。

张老栓把伤员小心地塞进去，让他半靠着土壁。又摸黑从自己床铺底下，拿出小半块黑乎乎、硬得像石头的杂面饼——这是他半天的口粮。

再提起墙角一个破葫芦，里面装着半葫芦清水。

他蹲在洞口，把饼子和水递进去。伤员似乎恢复了一点意识，艰难地抬起眼皮，眼神里充满了警惕和茫然。

“吃点，喝口。”张老栓的声音干巴巴的，没什么情绪， “别出声。白天不能出来。”

那伤员盯着他看了几秒，似乎判断出眼前这个满脸沟壑的老农没有恶意，才用颤抖的手接过饼和水，费力地咬了一口，又灌了口水。

张老栓不再说话，把柴捆仔细挪回原处，遮严实了洞口，又洒了点灰土掩饰痕迹。

然后他像没事人一样，躺回自己那张吱呀作响的破木板床上，睁着眼，听着屋后山风吹过树梢的呜咽，和柴垛后那极其轻微的、啃食干粮的声音。

第二天，村里果然不太平。

王老爷府上的几个家丁，陪着一个穿皮甲、挎腰刀的士兵，凶神恶煞地挨家挨户盘问，有没有见过生人，特别是带伤的。

说是有一股流匪在这一带活动，上头有令，见者立即禀报，隐藏不报者杀头。

到了张老栓家，兵爷用刀鞘胡乱拨拉着空荡荡的屋角，眼睛像鹰一样四处扫视。家丁踢了踢屋角的柴堆，柴火哗啦响了一下。张老栓的心提到了嗓子眼，垂着头，腰弯得更厉害，脸上是惯有的木然和惶恐。

“老东西，见过生人没有？”士兵喝问。

张老栓连连摇头，舌头像是打了结：“没…没有，军爷……咱这地方，鬼都不来……”

兵爷嫌恶地看了他一眼，又瞅了瞅那一贫如洗、一眼就能望到头的屋子，料定这老棺材瓤子也没胆子藏人，更没油水可榨，骂骂咧咧地走了。

听着脚步声远去，张老栓才缓缓直起腰，走到灶台边，借着添柴的机会，轻轻拍了拍柴堆。里面什么回应也没有，但他知道人还在。

往后的两天，张老栓白天照常下地，侍弄他那几分贫瘠的山田；晚上回来，会多煮一把野菜汤，盛出满满一碗，放在柴堆旁边一个不起眼的破瓦罐里。第二天早上，碗总是空了，洗得干干净净放回原处。饼子也时会少一点。

两人没有对话，一种无言的默契在潮湿的柴垛后和寂静的茅屋里流淌。

直到第三天夜里，伤员似乎恢复了些力气，能够自己稍微挪动了。张老栓搬开柴火，看见他正试图包扎自己手臂上一道较深的伤口，但动作笨拙。

张老栓沉默地走过去，从他手里接过洗干净的破布条——那是从他唯一一件还算完整的旧褂子上撕下来的。老人粗糙得像树皮的手，动作却出奇地稳当，仔细地帮他重新捆扎好。昏黄的油灯光晕下，一老一少，谁也没看谁。

包扎完了，张老栓默默坐回床边，准备吹灯。一直没怎么开过口的伤员，忽然在黑暗中低声说了一句：“……多谢老丈救命之恩。连累您了。”

张老栓动作顿了一下，没回头，半晌，才用他那干涩的、没什么起伏的乡音，木木地回道：

“谢个啥。”

停顿片刻，像是在思索，又像是在陈述一个再简单不过的事实，他补充道：

“你们打了王老爷的粮仓，却没抢俺家的鸡。”

黑暗中，柴垛后的呼吸声似乎凝滞了一瞬，随即，传来一声极轻的、像是释然，又像是苦涩的叹息。

张老栓吹熄了油灯。月光从破窗棂漏进来，照在他布满皱纹的、木然的脸上。他不懂什么大道理，不知道什么“问道途”，也分不清那女侠是真是假。

他只知道，王老爷的人来，准没好事，不是要粮就是要命，连他最后一只下蛋的母鸡上个月都被抢走了。

而柴垛后面那个年轻人，和他们一伙的，抢的是王老爷的粮仓，分给了像他一样或许永远去不了灰谷镇的穷人。

这就够了。

又过了几日，张老栓睡下之后，只听得堂下窸窸窣窣。

次日，他扒开柴堆，柴堆后那人已经走了，留下了一锭银子，约莫有一两之数。

这大致相当于他辛辛苦苦劳作一年的收入。

\vspace{2em}

最大胆，也最根本的支持，则是对改变命运之可能的直接拥抱。

那些纸张被摩挲得毛边、字迹因反复传抄而甚至出现了些错误的《养气篇》入门法诀，开始在深夜的油灯下、在矿洞交接班的间隙、在远离人烟的破庙里，被秘密地实践。

疲惫到极点的矿工，在浑身酸痛中勉强按照法诀指引，凝神感知。

十天，半月，毫无所感，只有更深的疲惫与沮丧。但总有一些人，凭借稍好的天赋或钢铁般的毅力，在某个万籁俱寂的深夜，忽然于丹田深处，捕捉到一丝微弱却无比真实、缓缓流转的暖意！

这丝暖流无法让他们力能扛鼎，却切实地缓解了积年的劳损酸痛，带来一种久违的、对自身身体与精神的微弱掌控感。

这种体验，如同在无尽黑暗中瞥见的一缕微光，一旦看见，便再也无法忘怀。

它成了最具感染力的“福音”，在绝对信任的兄弟、挚友间心照不宣地分享、探讨。

虽然范围依然受限，保密极为严格，但它像渗入千年旱地的涓涓细流，坚定、缓慢而无孔不入地扩展着“问道途”的根系。

星星之火，渐成燎原之势，其影响开始触及更高、也更复杂的层面。

一些在守旧联盟庞大体系中处于边缘、备受排挤、或是看清了联盟涸泽而渔本质的中小家族、不入流的小门派，开始暗中活动。

他们或许尚不敢公然易帜，但通过曲折隐晦的渠道，向问道途递出了橄榄枝：在联盟征调物资的清单上做手脚，截留部分品质最好的“孝敬”给吴刚的渠道；对自己辖地内那些深夜练气的“异动”睁一只眼闭一只眼，甚至暗中提供些许庇护。

更为大胆的人，会将族中那些天赋平平、不受重视，却对“问道途”理念心向往之的边缘子弟，以“游历”、“访友”之名悄悄送走，最终这些人往往出现在沈瑶或吴刚的队伍名录中；甚至有的人心甘情愿地做起了问道途的间谍。

散修中愿意投入游击队者，也终于不再仅限于空有一腔热血的青年；许多饱经世故、看透联盟将散修纯粹视为“耗材”本质的老江湖，在吴刚精明的利益折算和对散修心态入木三分的把握下，也做出了选择。

他们相信，在“问道途”这边，至少能被当作一个“人”来对待，流的血、出的力，能看到回报，更能赢得些许尊严。

沈瑶的游击队里，多了擅长在山林间追踪索迹的老猎户；吴刚的地下党中，加入了曾经走南闯北、熟悉各地方言与人情世故的“路路通”。

得益于此，一张无形却无处不在的耳目之网，在守旧联盟自以为严密的统治区内，悄然织就。

\vspace{2em}

情报战的态势，发生了根本性的、对联盟而言堪称灾难的逆转。

过去，守旧联盟的探子绞尽脑汁试图潜入云霞山、却往往有去无回；如今，在广大的乡村野地、散修市镇，一张无形而致密的耳目之网已然织就。

“青芒山一带，近日有残党啸聚，滋扰乡里，劫掠税赋。着令你部，于三日后卯时正，自黑水河西进，兵分两路，扫荡青芒山北麓及东侧河谷地带，务求全歼，以儆效尤！”

紫阳宗外门长老王猛将盖好印信的军令递给身旁身着文士衫、面相精明的师爷，沉声道：“速去誊抄三份，一份存档，一份送达黑石集驻军李统领，另一份快马传至青石城备案。此番必要肃清匪患，以振军威！”

“长老放心，属下即刻去办。”师爷恭敬地接过军令，躬身退下。

他回到自己那间堆满卷宗、弥漫着墨香与些许霉味的值房，掩上门。

脸上那副恭敬谨慎的神情如潮水般褪去，取而代之的是一种深沉的疲惫与一丝不易察觉的讥诮。

他展开军令，目光迅速扫过关键信息：三日后，卯时，黑水河西进，分两路，北麓与东侧河谷。

他先提起笔，在一张用于起草文书的草稿纸上，随意地记录了几个数字和简略方位。

这些符号混杂在其他无关的账目数字与地名缩写之中，即便被人看见，也只会当做寻常的工作笔记。

然后，他将这张草稿纸折好，塞进袖中一个特制的夹层。

接着，他才开始正式誊抄军令。笔走龙蛇，一丝不苟。

他写得很快，因为心中的紧迫感在催促。第一份存档，放入标有“乙未年七月剿匪”的卷宗匣内。第二份，装入专用的硬皮令筒，用火漆封好，盖上印鉴。

他唤来一名在门外候命的年轻传令兵。“将此令送至黑石集驻军李统领处，不得有误。”

“是！”传令兵接过令筒，转身便要走。

“且慢。”师爷叫住他，从案头拿起一个水囊递过去，“天热路远，带上这个。王统领脾气急，送达后莫要多言，速回复命即可。”

“多谢师爷！”传令兵感激地接过。

望着传令兵远去的背影，师爷坐回椅中，闭上眼，揉了揉眉心。

他又想起那个神秘的“行商”与自己“偶遇”时说的话：“先生大才，却屈居于此，替盘剥乡里、视人命如草芥之辈刀笔刑名，不觉得明珠暗投么？问道途所求，不过是让有才者得其用，有力者得其酬，这世道，不该是某些人生来就高高在上，而大多数人连抬头看一眼天空的机会都没有。”

当时他嗤之以鼻。但自己随后特意关注的一些关于联盟内部倾轧、以及云霞山内部分配制度的零星信息，却与他暗中观察到的一些迹象隐隐吻合。

更重要的是，对方也没有要求他做什么，只是说：“先生不妨冷眼旁观，看看这天下大势，究竟流向何方。选择全在先生身上。我等可不会逼迫先生。”

\vspace{2em}

命令还在路上，情报却已通过多条彼此不知晓的平行渠道，开始向地下网络的节点流淌。

黑石集，那个总是一脸和气的酒馆老板，正在柜台后擦拭酒杯。

一个风尘仆仆、贩运兽皮的货郎走进来，要了碗最便宜的粗茶，坐在角落默默喝着。货郎离开时，指尖似乎无意地在油腻的木桌上划过。老板过来收拾，目光扫过桌面，那里有几个茶渍水痕形成的、极其短暂的凌乱符号。

抹布一抹，符号消失。

但在打烊后，他进入后厨，从灶台旁一块松动的砖石后，取出一枚小小的玉简，将脑海中记忆的符号形状，以微弱的精神力刻画进去。

这枚玉简，会在明天清晨，由他那个恰好要去邻镇采买鲜货的侄儿带走。

几乎同时，在黑石集通往山外的岔路口，一个卖山货的农妇正在整理她的篮子，底层是沾着泥土的草药，上面盖着些野果。

一个看起来像是收药材的贩子蹲下来，挑拣着草药，低声讨价还价。最后生意没谈成，贩子摇摇头走了。

农妇也嘟囔着收拾篮子，但在无人注意时，她快速地从篮子底部的夹层里，抽出一根中空的细芦苇杆，里面卷着一条窄窄的布条，上面用炭笔写着歪斜的记号。

她不识字，但她知道要把这玩意传给下一个人。

这些零碎的情报，如同汇入溪流的无数水滴，通过吴刚精心编织、已愈发细密坚韧的网络，被迅速汇集、筛选、交叉验证。那个师爷的草稿符号、酒馆老板的玉简、农妇的布条……

它们从不同的社会阶层、不同的地理位置发出，沿着不同的安全路线流动，最终在某些绝对安全的隐蔽节点完成拼接和相互印证。

不到一天时间，一份完整、准确的敌情通报，已经摆在了沈瑶和吴刚的面前。

不止有清剿的时间、兵力、路线，甚至还包括了带队军官的性格特点、部分部队近期的士气状况。

\vspace{2em}

三日后，卯时。

黑水河西，紫阳宗外门弟子李统领骑着高大的龙血马，看着麾下三百余名修士和辅兵组成的队伍，在晨雾中列队完毕，刀枪映着微凉的晨光。

他踌躇满志，这次若能一举剿灭青芒山的“残党”，可是大功一件。

“出发！按计划，一队随我走北麓，二队走东侧河谷！仔细搜索，不得放过任何可疑痕迹！”王统领挥手下令。

大军开拔，脚步声、马蹄声、铠甲碰撞声打破了山林的寂静。

他们按照命令，仔细地搜遍了北麓的每一片林子，东侧河谷的每一个山洞。然而，除了发现几处早已废弃的临时营地痕迹、以及一些疑似有人活动过但已无从追查的细微线索外，一无所获。

连个“残党”的影子都没看到！

王统领的脸色从最初的踌躇满志，逐渐变得阴沉，最后涨成了猪肝色。

“怎么可能？！难道他们能未卜先知？！”

他愤怒地咆哮，心中却升起一股寒意。

这种一拳打在空处的感觉，比惨烈的战斗更让人憋屈和不安。

而就在李统领的大军在青芒山徒劳搜捕的同时，距离青芒山一百五十里外，一处隶属于丹鼎阁某附属家族的小型物资中转站，迎来了灭顶之灾。

这里存放着近期从附近几个矿点征收上来、尚未运走的低阶矿石和一批药材。因为原本驻扎于此的部分护卫被抽调去参与青芒山的“大行动”，守备力量比平日薄弱了三成不止。

留守的管家和几名护卫修士还在抱怨自己被留下看仓库是“没油水的差事”。

午时刚过，正是人最困乏的时候。数道如同鬼魅般的身影从仓库围墙的阴影中、从堆叠的货箱后悄无声息地闪现。

刀光如冷电，瞬间割断了哨兵的喉咙；沈瑶的身影如同融入阳光的幻影，直接出现在那名正在打盹的灵枢境护卫头领面前，在他惊骇睁眼的刹那，冰冷的锋刃已洞穿其护体灵光。

战斗短促而激烈，完全是一场不对等的屠杀。

负隅顽抗者被迅速格杀，其余人见势不妙，立即丢下武器投降。

沈瑶也没有管这帮家伙，只是让游击队员看着这群人，不让他们乱动。

迅速打开仓库，将里面堆积的矿石和药材能带走的带走；带不走的，连同仓库建筑，付之一炬。

冲天而起的黑烟，方圆数十里都清晰可见。

在撤离前，让队员将几袋便于携带的粮食，故意“遗落”在附近一个以贫苦闻名的村子外。

\vspace{2em}

广泛的群众基础和重新构建起的活跃的情报网络，使得游击战术的精髓得以淋漓尽致地发挥。

沈瑶手下精干的队伍，可以随时拆分成三四个更小的作战单元，如同水银泻地，同时袭扰数条补给线上的多个节点，焚毁粮草，破坏关键桥梁或运输法器，令敌军顾此失彼。

而沈瑶本人，则择机聚合精锐，执行高价值目标的斩首；那些尤其酷烈、贪婪、“没有团结价值”（吴刚语）的中低层军官，往往在某个酒醉的夜晚或巡视的路上，便被一道不知从何处而来的流光夺去性命。

这些行动，单次看来或许无法歼灭敌军成建制的主力，无法夺取一座城池，但其累积效应是恐怖而致命的。

守旧联盟漫长的补给线变得千疮百孔，从后方运往前线的物资损耗率急剧飙升，护送任务成了人人避之不及的死亡差事。

军队的士气，再难用简单的劫掠许诺和“必胜”口号来维持，厌战与恐惧的情绪在底层士兵和低级军官中如瘟疫般悄无声息地蔓延。

更让坐镇中军的陈炎焦头烂额的是，为了扑灭后方的“烽烟”，保护这些维系前线存在的“血管”，他不得不一次又一次地从围攻云霞山的主力中，分派精锐部队回援。

这些部队疲于奔命，穿梭在陌生的乡村山野，要么找不到神出鬼没的游击队，空耗钱粮士气；要么一脚踏入精心设计的伏击圈，损兵折将，狼狈不堪。

云霞山正面战场所承受的压力，就在这敌后持续不断的、精准的放血战术下，得到了切实而显著的缓解。山上的将士们发现，敌军的攻势间隔变长了，强度也有所减弱；而山下敌营中，因伤兵增多、补给不继而引发的抱怨与混乱，时有耳闻。

陈炎的军队自然也不是铁板一块：屡不建功、山下干耗许久，使得王丹和刘长靖对陈炎也失去了信心。两人对陈炎也颇有微词起来，守旧联盟内部也第一次出现了裂痕。

正是在这种背景下，陈炎在又一次接到后方某重要粮仓被焚、押运队全军覆没的急报后，终于暴怒而无奈地做出了战略调整。

他脸色铁青地打断了帐中关于下一次试探性进攻方案——当然也包括“究竟要不要打”——的争吵，嘶声道：“云霞山，不打了！传令，主力分批回撤！”

他不得不承认，继续将主力钉死在云霞山下，已是徒劳。

他计划逐步抽离主力部队，只在山下保持足够的兵力维持包围态势，防止山上守军大规模突围。

撤回的兵力，一部分用于加固几条核心补给线的防护，打造成“铁桶”通道；另一部分则组成数支机动精锐“黑旗营”，派往那些报告显示“匪患最烈”、“民心最浮”的区域，进行严厉无情的清洗与绥靖，意图用铁血手段扑灭所有反抗火苗，重新稳固后方。

与此同时，他试图以战利品分配优先权、未来地盘划分等空头许诺，拉拢那些在地方盘根错节的豪强家族，希望借助他们的“本土智慧”和人脉，来分辨良莠，协助维持基层。

然而，这剂猛药非但没能治病，反而成了最烈的毒药。

\vspace{2em}

那些在地方上经营数代、与守旧联盟中央嫡系本就存在利益纠葛与矛盾的中小家族，对陈炎这迟来的、充满功利色彩的“拉拢”看得一清二楚。

在一座小城，城主刘家接到了配合军队清剿境内“问道途暗桩”的严令。

刘家家主，一个面白无须、眼神精明的中年修士，恭敬地送走了杀气腾腾的联盟特使，转身回到密室，脸上的笑容瞬间消退。

“让我刘家去当这个恶人，替他紫阳宗趟雷？”他冷笑一声，对几位族老道，“问道途的人是那么好抓的？他们神出鬼没，在穷鬼里名声还好。我们真下死力去抓，抓得到几个不好说，这十里八乡的怨气，可全得记在我刘家头上。到时候，他紫阳宗拍拍屁股走了，我们还得在这片地上过日子！”

另一位族老捻着胡须，阴恻恻地补充：“正是。况且，谁能保证这问道途将来成不了气候？我看那叶牧谦，不是池中之物。现在把事做绝，将来何以自处？”

于是，刘家开始阳奉阴违起来。

他们派出的协查队伍，多是些老弱惫懒的家丁，每日在几个固定的、早就空无一人的废弃矿洞或山林外围转悠一圈，便算交了差。

偶尔上面催得急了，他们便去佃户里寻两个平日最不服管教、或是家里无力缴纳足额租子，做事还不老实的刺头，扣上“通匪”的帽子，五花大绑送去交差。

至于那些真正可能藏有伤员的村落，或者与吴刚网络有零星贸易往来的散修集市，刘家的人马总是“恰好”错过，或者“收到风声太迟”。

在更北边的平昌镇，地方豪强郝氏的做法则更显“高明”。

他们一方面大张旗鼓地配合联盟军队设立关卡，盘查行人，做出尽心竭力的姿态，甚至主动“贡献”出家族控制的几处偏远山庄，作为联盟军队的临时驻地。

但是另一方面，郝家族长却秘密吩咐心腹管家，通过绝对可靠的渠道，向吴刚的地下网络传递了一份经过筛选、不涉及核心利益的“清洗计划时间表”，同时附上一句：“迫于形势，虚应故事，万望海涵。贵方活动，请避开地图上划线两处。”

他们是在两头下注，用最小的代价，敷衍最大的凶险，同时为自己留一条后路。

这种“出工不出力”，抓替罪羊敷衍了事的现象，在非联盟嫡系的地方势力中成了普遍潜规则。

陈炎试图依靠的“本土秩序维护者”，首先维护的是他们自身的存续与利益，而非联盟那遥不可及且充满不确定性的“大业”。

这也是非常正常的。

哪个年代都不缺两头下注的骑墙派！

但这种骑墙派确实给问道途降低了巨量的压力！

\vspace{2em}

而那些被陈炎寄予厚望、授予全权的联盟嫡系机动精锐，则彻头彻尾地化为了噩梦的代名词。

这些“机动精锐”大多由紫阳宗失意弟子、依附家族的亡命徒、以及此前在正面战场或因剿匪不力受过处分、急于戴罪立功的军官组成，凶残暴戾，且被赋予“便宜行事”、“毋枉毋纵”的尚方宝剑。

他们面对的，也是陈炎和后方指挥部根本无法理解的现实困境：在问道途群众路线深耕已久的区域，“匪”与“民”的界限早已模糊不清，甚至浑然一体。

几乎家家户户都曾受过问道途小队或明或暗的恩惠——或许是深夜悄悄放在门前的半袋粮食，或许是家人重病时“恰好”路过游方郎中留下的药方，又或许只是在联盟税吏鞭打老人时，有陌生人站出来说了一句公道话。

几乎每个村落，都有年轻人偷偷修炼那本广为流传的《养气篇》，在体内生出第一缕微弱的、属于自己的元气。

对于这些手持利刃、背负军令、且本身就怀着抢掠发财欲望的“清剿”部队而言，分辨谁是“问道途”，成了不可能完成的任务。

部分人倒是能反应过来、悬崖勒马；但更多红了眼的人，认为分不清就不分了。

“传令麾下四王子，破城不须封刀匕！”

“宁可错杀一百，不可放过一人！”

“清洗”迅速、且必然地滑向了最野蛮、最血腥的深渊。

“黑旗营”开进村落之日，便是地狱洞开之时。

士兵们以搜捕“残党”、“收缴违禁功法”为名，如狼似虎地砸开每一扇门。

他们不在乎证据，只凭“眼神闪烁”“家中存粮稍多”“搜出写有不明字迹的纸片”等莫须有的理由，便将户主拖出，当着全家老小的面严刑拷打，逼问“同党”。

惨叫与哀求声响彻村落，但这只是前奏。

真正的灾难随之降临。

抢掠成了公开的行动。士兵们哄抢一切看得上眼的财物——不多的铜钱银角、稍微像样的衣物、甚至锅碗瓢盆。稍有反抗，便是刀剑加身。

村里的几个姑娘被拖进柴房，哭喊声很快被淫邪的笑声淹没。

一个老人仅仅是因为护着家里唯一的一只下蛋母鸡，便被当胸一刀刺倒。

最后，为了“以儆效尤”，也为了掩盖暴行，带队的校官下令点燃了几处被指认为“可能藏匿匪类”的房屋。

冲天的火光吞噬了屋舍，也吞噬了里面来不及逃出或不愿离开祖宅的老人。浓烟滚滚，混合着焦糊与血腥的气味，笼罩了这个曾经宁静的村落。

同样的场景，在多个被划为“重点清洗区”的村镇反复上演。

联盟军队带来的不是他们承诺的“秩序”，而是比传闻中任何“匪患”都更加无常、更加恐怖的噩梦。

在村民眼中，这帮疯子根本不是来平息动乱的官军，而是比盗匪更凶残百倍！

所过之处，十室九空，哀鸿遍野。

这种无差别的恐怖高压，在短期内确实取得了残酷的“效果”：不加区分的开杀，虽然大多数是错杀了，但也肯定有能歪打正着抓到几个真为问道途效力的。

一批忠诚而富有经验的联络员被捕后公开处决，首级被悬挂在城头示众。

沈瑶的游击队被迫化整为零，潜入更深的深山老林，活动几乎停滞，数次在针对性的围追堵截中险死还生，队伍减员严重。

吴刚的地下网络也遭受重创，许多线人被迫沉默或转移，情报传递变得异常艰难迟缓。

从表面战果看，陈炎的“清洗”似乎一度扼住了问道途敌后活动的咽喉，扑灭了若干“火焰”。但这恰恰是历史辩证法最无情的展现：高压，往往催生更猛烈、也更隐蔽的反抗。

表面的火焰被扑灭，灰烬之下，却是被烧得更透、更炽热、蔓延更广的草根与土壤！

粗暴野蛮的清洗，完成了问道途宣传队永远无法做到的、最彻底最深刻的动员。它将问道途这颗“杂草”的种子，用血与火作为肥料，更深、更牢地埋进了几乎每一个亲历者、乃至听闻者心底。

对于受过“官兵”劫掠的幸存者而言，当联盟士兵的刀砍向护着母鸡的老人、火把扔向祖屋时，所有的犹豫、恐惧和对“官府”残存的、惯性般的敬畏，都在那一刻燃烧殆尽了。

那个曾经只在夜间偷偷流传的“叶先生”和“问道途”，不再是一个模糊的符号或值得同情的“反贼”，而成了血海深仇中唯一的指望，活下去并寻求复仇的唯一可能路径。

甚至那些原本只是沉默忍受、内心隐隐偏向“秩序”——哪怕是不公的秩序——只求安稳度日的富农、小地主们，也在风暴中被席卷。

联盟军队并不总能准确区分“匪”和“民”——甚至更多的所谓“黑旗营官兵”已经不屑于去分辨了。

在抢红了眼的状态下，“通匪”的帽子可以扣给任何人，只为夺取其积蓄的财物。

当“秩序”的化身亲自演示了何谓无法无天的豺狼行径时，打破这“秩序”的“匪”，在百姓眼中便有了截然不同、甚至带上了悲壮与正义色彩的含义。

五帝三皇神圣事，骗了无涯过客；有多少风流人物？

盗跖庄蹻流誉后，更陈王奋起挥黄钺！

\vspace{2em}

民心的转变，不再是悄然的流淌，而是在绝望与暴政的熔炉中，完成了彻底、决绝的淬炼与逆转；也正是在这血腥的镇压与更顽强、更智慧的隐蔽反抗中，一种全新的、守旧联盟高层思维完全无法理解的战略格局，瓜熟蒂落，悄然成形。

云霞山，作为问道途不可撼动的战略支点、精神灯塔与正面战场，依然巍峨矗立，牢牢吸引并消耗着联盟剩余的注意力和部分兵力。

而广袤无垠的乡村野地、散修聚集的市镇、远离核心统治的中小城池——这些守旧联盟统治链条末梢、力量相对薄弱的区域——却在血腥的洗礼后，于无声处完成了人心的皈依，变成了问道途的海洋，革命的根据地初见雏形。

这里没有问道途的大军驻扎，没有飘扬的显赫旗帜，却有着无数双自愿为问道途守望的眼睛，无数张自愿为问道途传递消息的嘴，无数颗在黑夜中默默祈祷云霞山胜利、期待“叶先生”理念成真的心。

诉苦大会在深夜里秘密召开，红黑账本在暗地里严谨记录，互助性的、脱离守旧联盟掌控的自治组织在废墟上萌芽。

尽管守旧联盟的正规军依然控制着主要城池和交通要道，但他们惊恐地发现，一旦离开城墙的庇护，进入广大的乡村，他们就像滴入大海的墨点，被无尽的沉默与潜在的敌意所包围。

征粮变得异常困难，巡逻队时常遭遇冷箭或陷阱，甚至整队人马在熟悉的道路上诡异地“迷路”失踪。

联盟的军队可以凭借武力占领或摧毁任何一座他们看得见的城池，可以剿灭任何一支他们能追踪到的游击队，但他们永远无法占领每一寸土地，无法控制每一颗人心。

他们仿佛陷入了一片无形的沼泽，看似站在坚实的城池据点之上，实则举步维艰，因为承载他们的整个大地，已然悄悄改变了颜色。

农村包围城市——这一超越纯粹军事范畴、深植于社会土壤与人民意志的战略格局，在血与火的残酷浇灌下，在希望与绝望的反复淬炼中，终于第二次脱离纸面的构想，第一次走出历史的记载，而成为了内域大地上活生生的、呼吸着的现实。

上一次走向现实，还是叶牧谦的前世。

尽管敌后战场因残酷清洗而满目疮痍，付出了鲜血与牺牲的代价，但最深层的、决定战争最终走向的力量对比，已经发生了不可逆转的根本性变化。

陈炎的退兵与清洗，非但不是挽回颓势的妙手，反而成了加速这一历史进程的最后、也是最有力的推手。

他所面对的，早已不再是一个单纯的军事对手叶牧谦，而是一个从旧世界痛苦母体中孕育诞生、已深深扎根于亿万民众渴望之中的全新体系与浩荡潮流。

